% 本文件由 LaTeX 排版 Agent 根据有效 translation.json 自动生成。
% author=Miscellaneous; work_id=07158; title=宗徒宪章（卷八）

\chapter{宗徒宪章（卷八）}

% structure: p0001 source=h1 target=omitted reason=page-title-already-rendered-by-parent

\Emph{论神恩、圣职与教会法典。}\par

% structure: p0003 source=h2 target=section reason=relative-heading-rank; base=h2

\section{第一节 论神恩的多样性}

% structure: p0004 source=h3 target=subsection reason=relative-heading-rank; base=h2

\subsection{论行奇迹之能力为谁而显}

一、\Person{耶稣基督}，我们的天主和救主，将虔诚的伟大奥秘传授给我们，并召唤犹太人和外邦人，承认祂的父是独一真天主，正如祂在某个地方为那些相信之人的救恩而感恩时说：\BibleQuote{我将祢的名显示给世人，我完成了祢交给我的工作；}\BibleRef{若 17:6} 并为我们向祂的父说：\BibleQuote{圣父啊，世人虽不认识祢，我却认识祢；这些人也认识祢。} 祂有充分的理由，当我们因那由圣神从祂那里赐下的恩赐而得以圆满时，对我们众人说：\BibleQuote{信我的人必有这些标记相随：因我的名驱逐魔鬼，说新语言，手拿毒蛇，若喝了什么毒物，也决不受害；按手在病人身上，可使他们痊愈。}\BibleRef{谷 16:17} 这些恩赐最初赐给我们宗徒，当时我们即将向一切受造物宣讲福音，后来也必然赐予那些因我们而相信的人；并非为了行这些事的人的利益，而是为说服不信的人，好叫那言语未能说服的，被标记的能力所羞愧：因为标记不是为我们这些相信的人，而是为不信的人，无论是犹太人还是外邦人。因为我们驱逐魔鬼并无益处，而是为那些藉主的权能得以洁净的人；正如主自己在某处教训我们并显示说：\BibleQuote{不要因为魔鬼屈服于你们而欢喜，而要因你们的名字已登记在天上而欢喜。}\BibleRef{路 10:20} 既然前者是凭祂的能力做成，后者却凭我们的善志和勤奋，但显然还是靠祂的助佑。因此，并非每一位信徒都必须驱逐魔鬼、复活死人、说异语，只有那蒙恩得此恩赐的人才需要，为的是有利于不信者的救恩，他们常被羞愧，不是因世界的明证，而是因标记的能力；也就是说，那些堪当得救的人：因为所有不敬神的人并不被奇事所动；对此天主亲自作证，如同祂在律法中说：\BibleQuote{我要用异语和异唇对这百姓说话，但他们仍不相信。}\BibleRef{依 28:11} 因为埃及人虽然见过\Person{梅瑟}行了那么多标记和奇事，仍不相信天主；犹太人的群众也不相信基督，虽然祂在他们中间医治了各种疾病和灾殃。前者未被那变为活蛇的棍杖、那生麻风变白的子、那变为血的尼罗河所羞愧；后者未被那复明的盲人、那行走的瘸子、那复活的死人所羞愧。前者受\Person{雅乃斯}和\Person{杨布勒}抵挡，后者受\Person{亚纳斯}和\Person{盖法}抵挡。可见，标记并不能使所有人羞愧而相信，只有那些心地善良的人才能；为他们的缘故，天主也乐意，如同一位聪明的管家，指定奇迹藉以施行，并非靠人的能力，而是凭祂自己的旨意。我们说这些事，是为了叫那些领受了这种恩赐的人不要高抬自己，凌驾于未领受的人之上；我们指的是行奇迹的恩赐。因为，没有一个人相信了天主、藉着基督而未领受某种神恩：因为，从多神教的不敬中被救出，并藉着基督相信了天主父，这本身便是天主的恩赐。脱去犹太教的遮面纱，并相信：按天主的善意，祂的独生子——万世之前所生——在末世由童贞女所生，无人伴随，并作为人生活，却没有罪，完成了律法的一切义；又按天主的允许，那本是天主的圣言忍受了十字架，轻看了羞辱；祂受死、埋葬，并在三日之内复活；复活后，与宗徒们一起四十天，完成了全部法令，然后在他们眼前被接升到差祂来的天主父那里。凡相信这些事，不是随意和非理性，而是以判断和确信，便领受了天主的恩赐。同样，那从一切异端中被救出的人也如此。所以，行标记和奇迹的人不可判断那未得同样恩赐的信徒：因为天主藉着基督所赐的恩赐各不相同；有人领受这个恩赐，另一个人领受那个恩赐。也许有人有智慧的言语，另一个人有知识的言语；另一个能辨别神类；另一个能预知未来；另一个有教导的言语；另一个有坚忍；另一个有按律法的节制：因为连天主的仆人\Person{梅瑟}，在埃及行标记时，也没有高抬自己胜过同辈；当他被称为神时，也没有骄傲地藐视自己的先知\Person{亚郎}。\Person{若苏厄}，\Person{农}的儿子，在他之后的百姓领袖，虽然在攻打\Place{耶步斯}人时使太阳停在\Place{基贝红}之上，月亮停在\Place{阿雅隆谷}之上，因为白昼不够他们得胜，他也没有侮辱\Person{丕乃哈斯}或\Person{加肋布}。\Person{撒慕尔}行了那么多奇事，也没有轻视天主的爱人\Person{达味}：他们两人都是先知，一个是司祭长，另一个是君王。当以色列中只有七千人没有向巴耳屈膝时，\BibleRef{列上 19:18} 他们中间的\Person{厄里亚}和他的门徒\Person{厄里叟}是行奇迹者。然而\Person{厄里亚}没有轻视敬畏天主的管家\Person{敖巴狄雅}，后者没有行标记；\Person{厄里叟}也没有轻视自己在敌人面前恐惧的门徒。\EditorialNote{列王纪下第六章} 此外，智慧的\Person{达尼尔}两次从狮口中得救，以及那三个从火窑中被救出的孩子，也没有轻视其余的以色列同胞：因为他们知道，自己逃脱这些可怕的灾难不是凭自己的力气，而是凭天主的权能；他们既行奇迹，又从灾难中得救。因此，你们中任何人不可高抬自己反对弟兄，即使他是先知，或行奇迹者：因为如果不再有不信的人，一切标记的能力便从此多余。因为敬神出于个人的善志；行奇事却出于那藉我们行奇事者的能力：前者关乎我们自己，后者关乎行奇事的天主，原因已提过。因此，国王不可轻视手下官员，统治者不可轻视臣民。因为没有被统治的人，统治者便多余；没有官员，王国便无法存立。再者，主教不可高抬自己反对他的执事和司铎，司铎也不可反对人民：因为教会的存续在于彼此相依。因为主教和司铎对于人民是司祭，平信徒对于圣职人员是平信徒。成为基督徒在于我们自己的意愿；但成为宗徒、主教或任何其他职务，不在于我们自己的意愿，而在于天主安排，祂赐予恩赐。关于那些蒙恩得恩赐和尊位的人，就说到这里。\par

% structure: p0006 source=h3 target=subsection reason=relative-heading-rank; base=h2

\subsection{论不相称的主教与司铎。}

II. 其次，我们补充说：并非每个说预言的人都是圣洁的，也并非每个驱魔的人都是虔诚的；因为甚至贝敖尔的儿子\Person{巴郎}预言者也曾说预言，尽管他自己是不虔敬的；同样，那伪称为大司祭的\Person{盖法}也是如此。不仅如此，魔鬼和邪魔关于祂预言了许多事；然而尽管如此，他们心中却没有一点虔诚；因为他们因自愿的邪恶而被无知压制。因此，显然不虔敬的人即使说预言，也不能用他们的预言掩盖自己的不虔敬；那些驱魔的人也不会因魔鬼服从他们而成圣；因为他们只是互相欺骗，如同那些玩幼稚把戏取乐的人，并毁灭那些听从他们的人。因为一个邪恶的国王不再是国王，而是暴君；一个被无知或恶意压制的主教也不是主教，而是伪称的，不是由天主派遣，而是由人派遣，如同耶路撒冷的\Person{阿纳尼雅}和\Person{撒默雅}，以及巴比伦的假先知\Person{则德克雅}和\Person{阿希雅}。的确，预言者\Person{巴郎}曾用培敖耳的\Place{巴耳}使以色列人堕落，遭受了惩罚；\Person{盖法}最终自杀身亡；\Person{斯盖瓦}的儿子们试图驱魔，反被恶魔所伤，狼狈逃窜；以色列和犹大的国王们一旦成为不虔敬者，就遭受了各种惩罚。因此，显而易见，那些同样伪称的主教和司铎也逃不脱天主的审判。因为现在也要对他们说：\Quote{哎，你们这些藐视我名的司祭，我必将你们交付屠杀，如同我对待\Person{则德克雅}和\Person{阿希雅}，\Place{巴比伦}王将他们放在煎锅里煎了，}正如先知\Person{耶肋米亚}所说。我们说这些话，并非蔑视真正的预言，因为我们知道这些是由天主感动在圣洁之人身上成就的，而是为了制止虚荣之人的胆大妄为；并加上这一点：天主从这样的人身上收回祂的恩宠，因为\Quote{天主拒绝骄傲的人，却赐恩宠给谦逊的人。}在我们时代，\Person{息拉}和\Person{阿加波}曾说过预言；然而他们并没有使自己与宗徒同等，也没有越过自己的本分，尽管他们是天主所爱的。妇女也曾说预言。在古代，\Person{米黎盎}是梅瑟和亚郎的姐妹，\BibleRef{出 15:20}；之后有\Person{德波辣}\BibleRef{民 4:4}；再后有\Person{胡耳达}\BibleRef{列下 22:14}和\Person{友弟德}\EditorialNote{友弟德传 8}——前者在约史雅时代，后者在达理阿时代。主的母亲也曾说预言，她的亲戚\Person{依撒伯尔}和\Person{亚纳}也是如此；在我们时代，\Person{斐理伯}的女儿们\BibleRef{宗 21:9}；然而她们并没有因此傲视自己的丈夫，而是守住了自己的本分。因此，如果你们中间也有男人或女人获得某种恩赐，让他保持谦卑，好使天主喜悦他。因为祂说：\Quote{但我所垂顾的，是谦卑、安静、敬畏我言语的人。}\BibleRef{依 66:2}\par

% structure: p0008 source=h2 target=section reason=relative-heading-rank; base=h2

\section{第二部分 主教的选举与祝圣：主日礼仪形式}

% structure: p0009 source=h3 target=subsection reason=relative-heading-rank; base=h2

\subsection{论制定关于在教会中执行职务的规章之重要性。}

III. 我们现在已经完成了关于各种恩赐的第一部分论述，这些恩赐是天主按自己的旨意赐予人的；以及祂如何斥责那些试图说谎或被敌人之灵感动之人的行径；并且天主常常利用恶人来预言和行奇事。但现在我们的论述急转至主要部分，即教会事务的规章，好使你们——由我们按基督的命令祝圣的主教们——在从我们学习了这些规章之后，能遵照所交付的命令行事，知道听从我们的就是听从基督，听从基督的就是听从祂的天主和父\BibleRef{路 10:16}，愿光荣归于祂，直到永远。阿们。\par

% structure: p0011 source=h3 target=subsection reason=relative-heading-rank; base=h2

\subsection{论圣秩授予。}

IV. 因此，我们——现今聚集在此的主的十二位宗徒——将这些关于每一种教会形式的神圣规章交付给你们，与我们同在的还有蒙拣选的器皿、我们的同宗徒\Person{保禄}，主教\Person{雅各伯}，以及其余的司铎和七位执事。首先，我\Person{伯多禄}说：被祝圣的主教应当像我们已经全体所规定的，在一切事上无可指责，是经过拣选的人，由全体人民选出，当他被提名并获得认可后，让人民与在场的司铎团和主教们在主日聚集，并让他们表示同意。然后，首席主教应询问司铎团和人民，此人是否是他们所愿作他们首领的。如果他们表示同意，他再进一步问，此人是否从众人那里获得好见证，证明他配得上如此伟大光荣的权柄；他对天主的虔诚是否正直；他对人所行的正义是否被遵守；他是否妥善管理了自己的家务；他在生活中是否无可指责。如果全体会众真实地、不带偏见地作证他是这样的人，那么他们第三次——如同在天主审判者、基督面前，圣神也在场，以及所有圣善服役的神灵面前——再问一次，他是否真正配得上这个职务，以便\Quote{凭两个或三个见证的口供，一切事都可成立。}\BibleRef{玛 18:16}如果他们第三次同意他是配得的，就应要求他们全体投票；当全体自愿投票时，就应听取他们的意见。沉默之后，让一位首席主教与另外两位一起站在祭台旁，其余的主教和司铎默默祈祷，执事们将神圣的福音书展开放在被祝圣者的头上，并这样向天主说：——\par

% structure: p0013 source=h3 target=subsection reason=relative-heading-rank; base=h2

\subsection{为祝圣主教之祈祷格式。}

五、噢，伟大的存在，主全能的天主，唯祢无始，不受管辖；祢恒常存在，在世界之前就已存在；祢一无所缺，超越一切原因和肇始；唯祢是真天主，唯祢是智慧；唯祢是至高者；祢本性不可见；祢的知识无始；唯祢是善，无可比拟；祢在万物存在之前就已认识它们；祢知晓最隐秘的事；祢不可接近，无人在祢之上；祢是祢独生子、我们的天主和救主的天主和父；祢藉祂创造了整个世界；祢的眷顾照顾并看管一切；祢是仁慈的父和一切安慰的天主\BibleRef{格后 1:3}；祢居于至高天，却垂顾下界的事物；祢藉基督在肉身的来临设立了教会的规条；圣神藉祢的宗徒，并藉我们这些因祢恩宠在此的主教，为这规条作证；祢从起初就预定了司铎来治理祢的子民——首先是亚伯尔、舍特和厄诺士、哈诺客和诺厄、默基瑟德和约伯；祢任命了亚巴郎和其余圣祖，以及祢忠信的仆人梅瑟和亚郎、厄肋阿匝尔和丕乃哈斯；祢从他们中间拣选了在祢盟约之幕中的统治者和司铎；祢拣选了撒慕尔为司铎和先知；祢没有使祢的圣所缺少事奉者；祢因那些蒙祢拣选以得光荣的人而喜悦。求祢藉着我们，藉着祢的基督的中介，倾注祢自由圣神的德能，这德能已托付给祢的爱子耶稣基督；祂按祢的旨意将圣神赐给了祢永恒天主的圣宗徒。求祢，洞察人心的天主，因祢的圣名赐予祢这位仆，就是祢拣选为主教的那位，使他能牧养祢的圣羊群，向祢履行大司祭的职分，并夜以继日无可指摘地事奉祢；使他能平息祢的义怒，聚集那将得救之人的数目，并向祢献上祢圣教会的礼品。全能的主，求祢藉着祢的基督赐给他圣神的共融，使他能按祢的命令有权赦罪；按祢的命令分派份额；按祢赐给宗徒的权柄解除一切束缚；使他能以柔和与纯洁的心、以坚定、无可指摘、无可责备的意念令祢喜悦；向祢献上纯全无血的祭献，就是祢藉祢的基督所设立的新盟约的奥迹，作为甘饴的馨香，藉着祢的圣子耶稣基督、我们的天主和救主；藉着祂，在圣神内，愿光荣、尊威和钦崇归于祢，从今时直到永远，万世无疆。当他为此祈祷后，其余的司铎应加上\Quote{阿们}；全体百姓也应一同说\Quote{阿们}。祈祷后，一位主教应将祭品举在受祝圣者的手上；清晨时，应将他安置在他的宝座上，就是在其他主教中为他预备的地方；他们众人都要在主内与他行亲吻礼。在诵读法律和先知书、我们的书信、宗徒大事录和福音之后，受祝圣者应向教会致候说：\Quote{我们的主耶稣基督的恩宠、天主父的爱和圣神的共融，与你们众人同在。}众人应回答：\Quote{也与你的心灵同在。}这些话之后，他应向百姓讲劝勉的话；当他讲完训导的话后（我，伯多禄的兄弟安德肋，这样说），全体起立，执事应登上高处，宣告：\Quote{任何听道者、任何不信者，不要留下。}静默后，他应说：——\par

% structure: p0015 source=h3 target=subsection reason=relative-heading-rank; base=h2

\subsection{牛津抄本}

五、我们的主耶稣基督的天主和父，仁慈的父和一切安慰的天主，祢在一切发生之前就已知道；祢藉着祢恩宠的言语设立了教会的规条；祢从起初就预定了那义人的族系，就是从亚巴郎出来的统治者，并立他们为司铎，没有使祢的圣所缺少事奉者；祢从创立世界以来，因那些蒙祢拣选以得光荣的人而喜悦；如今求祢倾注祢自由圣神的德能，这德能藉着祢的爱子耶稣基督，祢已赐给了祢的圣宗徒，他们建立了教会以代替圣所，为祢的名得无尽的荣耀和赞美。噢，洞察万心者，求祢赐予祢这位仆人，就是祢拣选为祢主教圣职的那位，使他能向祢履行大司祭的职分，并夜以继日无可指摘地事奉祢；使他能不停地平息祢的义怒，并向祢献上祢圣教会的礼品；使他以大司祭的精神，按祢的诫命有权赦罪，按祢的谕令分派份额，按祢赐给宗徒的权柄解除一切束缚，并令祢喜悦，以柔和及纯洁的心献上甘饴馨香的祭品，藉着祢的子耶稣基督我们的主；愿光荣、权能和尊崇与圣神一起归于祢，从今时直到永远。阿们。\par

% structure: p0017 source=h3 target=subsection reason=relative-heading-rank; base=h2

\subsection{神圣礼仪，其中包含为慕道者的求恩祈祷}

VI. 你们慕道者，祈祷吧，并让所有信众在心意中为他们祈祷，说：主，求祢怜悯他们。并让执事为他们献上祈祷，说：让我们大家为慕道者向天主祈祷，使那位良善的、热爱人类者，仁慈地垂听他们的祈祷和恳求，并如此接纳他们的呼求，以助佑他们，赐给他们心中那些对他们有益愿望，并向他们启示祂基督的福音；赐予他们光照和悟性，教导他们认识天主，教给他们祂的诫命和规条，在他们心中植入祂纯洁而救人的敬畏，开启他们心中的耳朵，使他们昼夜操练于祂的法律；坚固他们的虔诚，将他们联合并纳入祂的圣羊群；恩赐他们重生的沐浴和不朽的衣袍，即真正的生命；并拯救他们脱离一切不虔敬，不给仇敌留下任何反对他们的余地；\Quote{并洁净他们所有肉体和精神的污秽，住在他们内，并在他们内行走，藉着祂的基督；祝福他们的出入，为了他们的益处安排他们的事务。}让我们仍热切地为他们献上我们的恳求，使他们因被接纳而获得过犯的宽恕，如此堪当神圣奥迹以及恒常与圣人的共融。起来，你们慕道者，为你们自己恳求天主的平安，藉着祂的基督，一个平安的日子，无罪的，并类似地为你一生的全部时间，以及你的基督徒的终结；一位慈悲而仁爱的天主；以及你们过犯的宽恕。将你们自己奉献给唯一不受生的天主，藉着祂的基督。低头，领受祝福。但正如我们先前所说，当执事逐一呼名时，让百姓说：主，求祢怜悯他；并让孩子们先说。当他们低头后，让新祝圣的主教以如下祝福祝福他们：全能的天主，不受生的、不可接近的，唯有祢是真实的天主，祢的基督、祢的独生子的天主和父；施慰者的天主，并普世的主；祢藉着基督指派祢的门徒为教导虔敬的导师；求祢现在也垂顾祢的仆人们，他们正在接受祢基督福音的教诲，并\Quote{赐给他们一颗新心，在他们内里更新正直的精神，}使他们全心、甘愿地认识并承行祢的旨意。恩赐他们圣洁的接纳，将他们联合于祢的圣教会，使他们成为祢神圣奥迹的分享者，藉着基督，祂是我们的希望，并为他们而死；藉着祂，在圣神内，愿光荣和朝拜归于祢，直到永远。阿们。此后，让执事说：慕道者，平安出去。他们出去后，让他说：你们附魔者，被不洁之灵折磨的，祈祷吧，并让我们大家热切为他们祈祷，使热爱人类的天主藉着基督斥责不洁和邪恶的灵，并将祂的恳求者从仇敌的辖制下解救出来。愿祂曾斥责那群魔鬼和魔鬼、邪恶的首领，\BibleRef{谷 5:9}现在也斥责这些背弃虔敬者，并将祂自己的化工从他那权力下解救出来，并洁净祂以极大智慧所造的那些受造物。让我们仍为他们热切祈祷。拯救他们，天主，并以祢的大能提升他们。低头，你们附魔者，领受祝福。并让主教加上这个祈祷，说：—\par

% structure: p0019 source=h3 target=subsection reason=relative-heading-rank; base=h2

\subsection{为附魔者}

VII. 祢，曾捆绑壮士，夺取了他屋内的一切，祢赐给我们权柄践踏蛇和蝎子，并一切仇敌的能力；\BibleRef{玛 12:29}祢将那蛇，杀人的凶手，像给孩子的麻雀一样捆绑给我们，万物都畏惧，在祢能力面前战栗；祢曾如同闪电般将他从天上打到地上，不是地方的坠落，而是因他自愿的邪恶倾向，从荣耀到羞辱；祢的注视干涸深渊，祢的威吓融化山岭，祢的真理永远长存；婴孩赞美祢，吃奶的婴儿祝福祢；天使歌颂祢，朝拜祢；祢注视大地，大地战栗；祢触摸山岭，山岭冒烟；祢威胁大海，使之干涸，使所有河流变为旷野，云彩是祢脚下的尘土；祢步行海上如履平地；祢，独生天主，伟大圣父的圣子，斥责这些邪恶的灵，并将祢手的化工从敌对之灵的权力下解救出来。因为光荣、尊荣和朝拜归于祢，并藉着祢归于祢的父，在圣神内，直到永远。阿们。并让执事说：出去，你们附魔者。他们之后，让他大声呼喊：你们即将受光照者，祈祷吧。让我们所有信众热切为他们祈祷，使主恩赐他们，在被引入基督的死亡后，能与祂一同复活，成为祂王国的分享者，并被接纳进入祂奥迹的共融；将他们联合于、并数算在祂圣教会中那些得救的人中。拯救他们，并以祢的恩宠提升他们。并且，藉着祂的基督被印封归于天主后，让他们低头，从主教领受这个祝福：—\par

% structure: p0021 source=h3 target=subsection reason=relative-heading-rank; base=h2

\subsection{为受洗者}

八、祢曾藉祢的圣先知对那受归化的人说：\Quote{洗涤自洁，} \BibleRef{依 1:16} 并且藉着基督设立了属神的再生。求祢如今也垂顾这些受洗者，祝福他们，圣化他们，预备他们，使他们堪当祢属神的恩赐，以及祢属神奥迹的真正义子名分，得以与那些藉着我们的救主基督得救的人一同聚集。愿光荣、尊崇与钦崇藉着祂，在圣神内，归于祢，直到永远。阿们。执事说：\Quote{候洗者，请出去。}此后他宣告：\Quote{补赎者，请祈祷。让我们为补赎状态中的弟兄们全心祈祷，愿那怜爱人的天主指示他们悔改的道路，收纳他们的归回与告解，迅速将撒殚践踏在他们脚下，\BibleRef{罗 16:20} 救他们脱离魔鬼的罗网和恶灵的虐待，释放他们免受一切非法言语、一切荒唐行为和邪恶思想的束缚；宽恕他们所有故意的与非故意的过犯，涂抹那反对他们的手书，\BibleRef{哥 2:13} 将他们写在生命册上；\BibleRef{斐 4:3} 洁净他们身灵的一切污秽，\BibleRef{格后 7:1} 恢复他们并将他们重新联合到祂的圣羊群里。因为祂知道我们的形体。谁能夸耀自己有洁净的心呢？谁能坦然说自己纯洁无罪？\BibleRef{箴 20:9} 因为我们都在可责之列。让我们更为他们恳切祈祷，因为在天上一个罪人悔改，也有喜乐。\BibleRef{路 15:7} 好使他们远离一切恶行，转向一切善行；愿那爱人的天主迅速收纳他们的祈求，将祂救恩的喜乐归还他们，以祂自由的神坚固他们；使他们不再动摇，并被接纳与祂至圣的事物共融，成为祂神圣奥迹的分享者，以堪当祂的义子名分，获得永生。让我们都为他们恳切地说：主，求祢怜悯他们。天主，求祢拯救他们，以祢的慈悲使他们起来。请起来，藉着基督向天主低头，领受祝福。}然后主教加念此祷文：\newline{}\par

% structure: p0023 source=h3 target=subsection reason=relative-heading-rank; base=h2

\subsection{覆手；为补赎者的祈祷}

九、全能永生的天主，整个世界的主，万物的创造者与治理者，祢藉着基督将人展示为世界的装饰，并赐予他自然铭刻与书写的法律，使他作为理性的受造物按法律生活；当他犯罪时，祢赐予他祢的良善作为他悔改的保证。求祢垂顾这些向祢屈身灵与颈项的人；因为祢不喜悦罪人之死，却喜悦他悔改，转离他的恶道而生活。祢曾收纳尼尼微人的悔改，祢愿所有人都得救，并认识真理；\BibleRef{弟前 2:4} 祢曾以慈父心肠因那挥霍家产的儿子 \EditorialNote{路加福音第十五章} 的悔改而接纳了他；求祢如今接纳祢这些恳求者的悔改：因为没有人不犯罪；因为\Quote{主，祢若究察罪孽，谁能站得住呢？因为祢有赦免之恩。}求祢恢复他们进入祢的圣教会，回到他们原先的尊严与荣耀，藉着我们的天主和救主基督，愿光荣与钦崇藉着祂，在圣神内，归于祢，直到永远。阿们。然后执事说：\Quote{补赎者，请离开。}他又说：\Quote{那不该上前的人，谁也不要靠近。我们众信徒，让我们屈膝；让我们众人藉着基督恳求天主；让我们藉着基督热切祈求天主。}\newline{}\par

% structure: p0025 source=h3 target=subsection reason=relative-heading-rank; base=h2

\subsection{为信友的求祷}

十．让我们为世界的和平与幸福安定，以及各圣教会祈祷；愿全世界之天主赐予我们祂永久的和平，就是那不能从我们手中夺去的和平；愿祂保守我们圆满地实践那合乎虔诚的德行。让我们为从地极到地极的圣而公宗徒教会祈祷；愿天主保守它坚固不移，免于此生的波浪，直到世界终结，如同建立在磐石上；也为本地的圣堂区祈祷，愿全世界的主恩赐我们不停地追随祂天上的希望，并不懈地向祂偿还我们祈祷的债务。让我们为所有普天下正确分施祢真理之言的整个主教职祈祷。也为我们的雅各伯主教及其堂区祈祷；让我们为我们的克莱孟主教及其堂区祈祷；让我们为我们的埃乌奥迪主教及其堂区祈祷；让我们为我们的安尼亚诺主教及其堂区祈祷：愿慈悲的天主恩赐他们健康、尊荣和长寿，继续在祂的圣教会中，并赐予他们在虔诚和正义中可敬的晚年。让我们为我们的长老祈祷，愿主拯救他们脱离一切无理和邪恶的行为，恩赐他们健康且尊荣的长老职。让我们为所有在基督内的执事及服务人员祈祷，愿主赐给他们无可指责的职分。让我们为读经员、歌咏员、贞女、寡妇和孤儿祈祷。让我们为已婚者和生育者祈祷，愿主怜悯他们众人。让我们为圣洁生活的阉人祈祷。让我们为在节制和虔诚状态中的人祈祷。让我们为在圣教会内结果实并向穷人施舍的人祈祷。让我们为向主我们的天主献祭并奉献供物的人祈祷，愿天主，一切美善的泉源，以祂天上的恩赐回报他们，\BibleQuote{给他们今世百倍，并在来世永生；}\BibleRef{玛 19:29}并以永恒的事物代替他们短暂的事物，以天上的事物代替地上的事物。让我们为新受光照的弟兄祈祷，愿主坚强并坚固他们。让我们为患病受试炼的弟兄祈祷，愿主解救他们脱离一切疾病和一切病患，并使他们痊愈回到祂的圣教会。让我们为水陆旅行者祈祷。让我们为在矿场、流放、监狱和锁链中因主之名的人祈祷。让我们为受苦役折磨的人祈祷。让我们为我们的仇敌和恼恨我们的人祈祷。让我们为因主之名迫害我们的人祈祷，愿主平息他们的愤怒，驱散他们对我们怀有的怒气。让我们为教外和迷失道路的人祈祷，愿主使他们归化。让我们纪念教会中的婴孩，愿主使他们在敬畏中臻于完美，并在年龄上达到成熟。让我们彼此代祷，愿主以祂的恩宠保守我们直到终末，拯救我们脱离那邪恶者，以及一切作恶之徒的丑闻，并保守我们直到祂的天国。让我们为每一个基督信徒的灵魂祈祷。天主，求祢以祢的仁慈拯救我们，扶持我们。让我们起来，恳切祈祷，并藉着基督，将我们自己及彼此交托于永生之天主。然后让大司祭添加这篇祈祷，并说：——\par

% structure: p0027 source=h3 target=subsection reason=relative-heading-rank; base=h2

\subsection{为信友的祈祷格式。}

十一. 全能的主，至高者，居住在高天的圣者，安息于圣者之中，无始，唯一的主宰，祢曾藉基督赐给我们真知的宣讲，使我们认识祢的光荣和祢的名，这是祂向我们启示，为使我们理解的，求祢现在也藉着他垂顾这祢的羊群，拯救它脱离一切无知和邪恶的行为，并恩赐我们以真诚敬畏祢，以深情爱慕祢，并对祢的光荣怀有应有的崇敬。求祢恩待他们，仁慈对待他们，当他们向祢祈祷时垂听他们；并保守他们，使他们坚定、无可指责、无可责难，使他们身灵圣洁，没有瑕疵、皱褶或任何类似之物；而是完整无缺，没有一人有缺陷或不完美。我们的支柱，全能的天主，祢不偏待人，求祢作这祢子民的助手，祢曾用祢基督的宝血救赎了他们；作他们的保护者、帮助者、供给者、守护者，他们坚强的防御墙，他们的堡垒和安全。因为\Quote{没有人能从祢手中夺去：} \BibleRef{若 10:29} 因为没有别的神像祢；因为我们信赖祢。\Quote{请以真理圣化他们：因为祢的话是真理。} \BibleRef{若 17:17} 祢不徇私情，无人能欺骗祢，求祢拯救他们脱离各种疾病、各种病症、各种过犯、各种伤害和欺诈，\Quote{在敌人的恐惧前，从白天飞来的箭，从黑暗中行走的灾祸；} 并赐予他们那永生，即在祢的唯一圣子、我们的天主和救主基督内的永生，愿藉着他，在圣神内，光荣与朝拜归于祢，从今直到永远，世世无穷。阿们。此后，让执事说：\Quote{请留心。} 让主教向教会致候，并说：\Quote{天主的平安与你们众人同在。} 让群众回答：\Quote{也与你的心灵同在。} 让执事向众人说：\Quote{请以神圣之吻彼此致候。} 让圣职人员向主教致候，平信徒中的男人向男人致候，女人向女人致候。让儿童站在读经台旁；让另一位执事站在他们旁边，以免他们失序。让其他执事巡视并看守男人和女人，以免发生骚动，无人点头、耳语或打盹；让执事站在男人们的门口，副执事站在女人们的门口，以免有人出去，或门被打开，即使是为了一个信友，在奉献礼的时候。但要有一位副执事取水为司铎洗手，这象征那些献身于天主的灵魂的纯洁。\par

% structure: p0029 source=h3 target=subsection reason=relative-heading-rank; base=h2

\subsection{若望的兄弟、载伯德的儿子雅各伯的宪章。}

十二、我——\Person{雅各伯}，\Person{若望}的兄弟，\Person{载伯德}的儿子，说：执事应立刻说：\Quote{望教者一个也不可留，听道者一个也不可留，不信者一个也不可留，异端者一个也不可留。你们这些祈祷了以上祷词的人，请离去。母亲们领回自己的孩子；谁也不可对任何人怀怨；谁也不可虚伪而来；我们要怀着敬畏和战栗，直立主前，预备奉献。} 这事完成后，执事应将礼品带到祭台前的主教那里；司铎应站在他的左右，如同门徒站在师傅面前一样。但两位执事应站在祭台两旁，手拿由薄皮或孔雀羽毛或细布制成的扇子，悄悄赶走飞动的虫豸，使它们不得靠近杯爵。因此，大司祭应与司铎们一同独自祈祷；他应穿上发光的衣服，站在祭台前，用手在额上划十字圣号，并说：\Quote{全能天主的恩宠，我们的主耶稣基督的爱，以及圣神的共融，与你们众人同在。} 全体应同声说：\Quote{也与你的心灵同在。} 大司祭：\Quote{请举起你们的心神。} 全体百姓：\Quote{我们向主举起心神。} 大司祭：\Quote{让我们感谢主。} 全体百姓：\Quote{这是理所当然的。} 然后大司祭说：\Quote{在万有之前向祢——真天主，在万有之前存在的天主，歌唱赞美诗，是极合宜且正当的；\InnerQuote{从祂，天上地下的整个家族都得名；}\BibleRef{弗 3:15} 祢是惟一的未生者、无始者、无统治者、无主人者；一无所缺；一切美善的施予者；超越一切原因和受造者；恒常不变；万物从祢而来，如同从它们本原而来。因为祢是永恒的知识、永久的看见、未生的听、未受教的知识、按本性为首、存在的尺度、超乎一切数目；祢藉祢的独生子从无中创造了万有，但祢在万世之前以祢的旨意、能力和美善，无需任何工具生了祂——独生子、天主圣言、生活的智慧、\InnerQuote{一切受造的首生者、祢伟大计谋的天使、} 祢的大司祭，也是一切理智与可感本性的君王与主，祂在万有之前，万物藉祂而存在。因为祢，永恒的天主，藉祂创造了万有，并藉祂使祢相称的眷顾运行宇宙；因祢藉以赐予存在的同一者，也赐予了福祉：祢独生子的天主和父，祢藉祂在万有之前创造了革鲁宾和色辣芬、永世\OriginalTerm{永世}{æons}和万军、异能者和掌权者、宰制者和宝座者、总领天使和天使；在这一切之后，又藉祂创造了这个可见的世界及其中的一切。因为祢是那以穹苍铺张诸天、\InnerQuote{展开诸天如同帐幕幔子，} 以祢的旨意将地立于虚无之上者；祢固定了穹苍，预备了黑夜和白天；从祢的宝库发出光，在光离开时带来黑暗，使世上活动的众生得以休息；祢设定天上的太阳管辖白天，月亮管辖黑夜，并在天上刻上星群合唱队以赞美祢荣耀的威严；祢造水供饮用和洗涤，我们赖以呼吸的空气供呼吸和发声，藉着击打空气的舌头和与之合作的听觉，使言语被接收并落在其上时能感知；祢造火在黑暗中安慰我们，供给我们的缺乏，使我们得以取暖和照明；祢将大海与陆地分开，使前者可航行，后者可步行，并以大小生物充满前者，以驯服和野生的同样生物充满后者；以各种植物装饰它，以香草点缀它，以花卉美化它，以种籽使它富饶；祢设立了深渊，在四周造了巨大的凹地，容纳堆聚的咸水大海，\EditorialNote{约伯传第三十八章} 然而祢用最微小沙粒的屏障四面围住它们；\BibleRef{耶 5:22} 祢时而用风将其掀起至山高，时而平静如镜；时而以风暴激怒它，时而以宁静平息它，使航海者航行容易；祢以河流环绕祢藉基督所造的这个世界，以水流浇灌它，以永不枯竭的泉源湿润它，以山脉围绕它，使之稳固不移。因为祢已充满祢的世界，以芬芳和医治的香草、众多不同的生物——强壮的与软弱的，为食物和劳力，驯服和野生；以爬行物的声音、各类飞鸟的声响；以年岁的循环、月日的数目、季节的次序、雨云的路径，为果实的生产和生物的供养。祢也设立了风的位置，它们按祢的命令吹动，以及众多植物和香草。祢不仅创造了世界本身，还造了人作为世界的公民，展现他为世界的装饰；因祢曾对祢的智慧说：\BibleQuote{让我们照我们的肖像，按我们的模样造人；使他们管理海里的鱼、天空的飞鸟。}\BibleRef{创 1:26} 因此祢也造了他，有不能死的灵魂和可朽坏的身体——前者出自虚无，后者出于四元素——并赐给他灵魂以理性知识，辨别虔诚和不虔，观察正义和不义；赐给他身体以五种感官和行进运动。因为祢，全能的天主，藉着祢的基督，在伊甸，\BibleRef{创 2:8} 东方栽植了一个乐园，以所有可食的植物装饰了它，并把他领进去，如同进入丰盛的筵席。当祢造他时，祢在他内植入了一条法律，使他能在自己家中和内在拥有神圣知识的种子；当祢把他领进喜乐的乐园时，祢赐他享受万物的特权，只禁止尝一棵树，期待更大的福乐；这样，如果他遵守那条诫命，就可得到其赏报——不死。但是当他忽略了那条诫命，在蛇的诱惑和妻子的劝说下尝了禁果时，祢公义地把他逐出了乐园。}\par

然而祢因自己的良善并没有忽略他，也没有任凭他完全灭亡，因为他原是祢的受造物；祢反而将整个造化置于他之下，赐给他自由，使他能以自己的汗水和劳力获取食物，同时使地上的一切果实发芽、生长并成熟。但当他暂时安睡后，祢以誓言呼召他重新复原，解开了死亡的束缚，应许他复活后的生命。不仅如此，当祢使他的后裔增加到无数时，那些与祢同住的，祢荣耀了他们；那些背离祢的，祢惩罚了他们。祢接受了亚伯（\EditorialNote{创世纪第四章}）的祭献，视之为圣者的祭献，却拒绝了杀害兄弟的加音的供物，视之为可憎的恶人。此外，祢悦纳了舍特和厄诺士\BibleRef{德 49:16}，并提携了哈诺客；因为祢是人的创造者、生命的赐予者、匮乏的供应者、法律的颁布者、遵守者的酬报者、违犯者的惩罚者。祢因不敬神者的众多，使大洪水临于世界，却藉方舟拯救了义人诺厄\BibleRef{伯前 3:20}，带着八条灵魂，成为前代世代的终结和未来世代的开始。祢点燃了可怕的烈火，攻击索多玛的五座城，\Quote{使肥沃之地变为盐泽，因其中居民的邪恶}，却从火焰中救出了圣者罗特。祢是那将亚巴郎从他祖先的不敬中解救出来、立他为世界的继承者、并将祢的基督启示给他的一位。祢预先指定默基瑟德为大司祭，侍奉祢的敬拜。祢使祢忍耐的仆人约伯战胜那恶之守护者的毒蛇。祢使依撒格成为应许之子，使雅各伯成为十二位儿子的父亲，并将他的后裔增多如海沙，带他带着七十五人进入埃及。主啊，祢没有忽略若瑟，反而因他为祢而守洁的功德，赐给他治理埃及人的权柄。主啊，祢没有忽略因埃及人压迫而受苦的希伯来人，因祢对他们的祖先的应许；祢拯救了他们，惩罚了埃及人。当人类败坏了自然律，有时将造化视为偶然的结果，有时过份尊崇它，将其等同于宇宙之神，祢并没有任凭他们迷失，反而兴起祢的圣仆梅瑟，藉他颁布了成文的法律以辅助自然律，并显明造化是祢的工程，驱散了多神论的谬误。祢以司祭职装饰了亚郎及其后裔，在希伯来人犯罪时惩罚他们，在他们回归祢时又接纳他们。祢以十灾难审判了埃及人，分开大海，带领以色列人通过，淹没并毁灭了追赶他们的埃及人。祢以木使苦水变甜；祢从磐石中引出水来；祢从天上降下玛纳，并以鹌鹑作肉食从空中降下；祢以火柱夜间光照他们，以云柱白日遮蔽他们免受炎热；祢宣告若苏厄为军队的统帅，藉他推翻了客纳罕的七个民族；祢分开约但河，使厄堂的河水干涸；祢不用器械或人手推倒了城墙。为这一切，全能的主，愿荣耀归于祢。无数的天使、总领天使、上座者、宰制者、率领者、掌权者、异能者，祢的永恒军队，都朝拜祢。革鲁宾和六翼色辣芬，用两翅遮盖脚，两翅遮盖头，两翅飞翔，与千千万万的总领天使、万万的天使，不停息地，以恒常而洪亮的声音说，愿全体百姓与他们同说：\Quote{圣，圣，圣，万军的上主，天地充满祂的荣耀：愿祂永受赞美。阿们。}\BibleRef{依 6:3}此后，大司祭应说：因为祢确实是圣的，至圣的，至高者，永远至尊。祢的独生子、我们的主和天主、耶稣基督也是圣的，祂在祢各种创造和合宜的眷顾中，事事服事祂的天主和父，并没有忽略迷失的人类。但在自然律之后，在成文法律的劝诫之后，在先知的斥责和天使的管理之后，当人类既败坏成文法律又败坏自然律，并抛诸脑后忘记了洪水、索多玛的焚烧、埃及的灾难和客纳罕居民的屠戮，几乎要以无与伦比的方式完全灭亡时，祂乐意因祢的善意而成为人——祂原是人的创造者；成为受法律约束者——祂原是立法者；成为祭品——祂原是大司祭；成为羔羊——祂原是牧者。祂平息了祢，祂的天主和父，使祢与世界和好，使众人免于将来的义怒，并由童贞女所生，取得肉身，祂是天主圣言，爱子，一切受造物的首生者，并依照关于祂所预言的那些预言（由祂自己预言），是达味和亚巴郎的后裔，出自犹大支派。祂在童贞女的胎中形成——祂是形成一切出生世人者；祂取得肉身——祂原无肉身；祂在时间之前受生，而在时间中诞生；祂圣洁地生活，并依照法律施教；祂驱除人间的各种疾病和灾殃，在百姓中行奇迹异事；祂分享饮食和睡眠——祂是一切需要食物者的养育者，\Quote{以祂的美善充满一切生灵}；\Quote{祂向不认识祂名的人启示了祂的名}\BibleRef{若 17:6}；祂驱除了无知；祂复兴了虔诚，满全了祢的旨意；祂完成了祢交托祂的工作；当祂安排好这一切后，祂被不敬神者的手——假称为大司祭和司祭者以及悖逆百姓的手——所擒拿，被那如同沉疴般拥有邪恶者所出卖；祂从他们手中受了许多苦，并因祢的许可忍受了各种羞辱；祂被解送到总督比拉多面前，审判者受审判，救主被定罪；那不能受苦者被钉在十字架上，那本性不死者死了，那赐予生命者被埋葬，为要释放那些祂因他们而来者脱离苦难和死亡，并折断魔鬼的锁链，拯救人类脱离他的欺骗。第三日祂从死者中复活；与祂的门徒共处四十天后，祂被提升天，坐在祢——祂的天主和父——的右边。\par

我们记念祂为我们所承受的一切，为此我们感谢祢，全能的天主，虽不配我们的感谢，却尽我们所能，并履行祂的建定：\Quote{因为祂被出卖的那一夜，祂拿起饼来}\BibleRef{格前 11:23}，在祂圣洁无玷的手中，举目仰望祢、祂的天主和父，\Quote{擘开，递给祂的门徒说：这是新约的奥迹：你们拿去吃。这是我的身体，为多人擘开，以得罪的赦免。}同样，祂\Quote{也拿起杯来}，用酒和水调和，祝圣后，递给他们说：\Quote{你们都喝这个，因为这是我的血，为多人流出，以得罪的赦免；你们要这样做，来纪念我。因为你们每次吃这饼，喝这杯，就是宣告主的死亡，直到我来。}因此，我们记念祂的苦难、死亡、从死者中复活、升天，以及祂未来的第二次来临，那时祂要在荣耀和权能中来临，审判生者死者，按照各人的行为予以回报，我们向祢，我们的君王和天主，按照祂的建定，奉献这饼和这杯，藉着祂感谢祢，因为祢使我们配得站立在祢面前，向祢祭献；我们恳求祢，仁慈垂视这些摆在祢面前的供物，天主，祢并不需要我们的任何奉献。请接纳它们，以光荣祢的基督，并派遣祢的圣神，主耶稣苦难的见证者，降临在这祭献上，使祂显示这饼是祢基督的身体，这杯是祢基督的血，好让领受者得以坚强于虔敬，获得罪的赦免，脱离魔鬼及其欺骗，充满圣神，堪当承受祢的基督，并在祢与他们和好时获得永生，全能的主。主，我们还要为祢的圣教会祈祷，她从世界的一端到另一端，是祢用祢基督的宝血所赎买的，求祢保护她坚定不移，不受干扰，直到世界终结；并为主教职，凡正确分解真理之言的。我们还为我自己祈祷，我一无所有，却向祢奉献，并为整个司铎团、执事及所有圣职人员祈祷，求祢使他们充满智慧，以圣神充满他们。主，我们还要\Quote{为君王和一切有权位者}\BibleRef{弟前 2:2}，并为全军祈祷，使他们与我们和平相处，好让我们一生在安宁与同心合意中度过，藉着耶稣基督——我们的希望——光荣祢。我们还为所有从创世以来蒙祢悦纳的圣者奉献——圣祖、先知、义人、宗徒、殉道者、宣信者、主教、司铎、执事、副执事、读经者、歌咏者、贞女、寡妇及平信徒，以及所有祢知道名字的人。我们还为这百姓奉献，求祢使他们成为\Quote{君尊的司祭民族、圣洁的国民}\BibleRef{伯前 2:9}，以赞美祢的基督；并为守贞和纯洁的人；为教会的寡妇；为尊贵婚姻和生育的人；为祢子民的婴孩，求祢不容我们任何人\Quote{成为可弃绝的}。我们还恳求祢为这城及其居民；为患病者；为在苦役中的人；为被放逐者；为在监牢里的人；为在水上或陆上旅行的人；祢是众人的帮助和支援，求祢作他们的支持。我们还恳求祢为那些因祢的名恨我们并迫害我们的人；为那些在外迷路的人；求祢使他们转回良善，平息他们的愤怒。我们还恳求祢为教会的慕道者，为被仇敌困扰的人，并为我们的忏悔者弟兄们，求祢使前者在信德上成全，使后者脱离恶者的势力，并接纳后者的忏悔，宽恕他们和我们的过犯。我们还为良好的气候和果实的丰收奉献，好让我们永远分享从祢而来的美物，不停赞美祢，赐食物给一切有血肉的。我们还恳求祢为那些因正当理由缺席的人，求祢使我们众人保守虔敬，并在祢基督——一切可感与理智本性的天主——的王国中聚集我们，我们的君王，求祢使我们坚定不移、无可指责、无可非难：因为一切荣耀、敬拜、感谢、尊荣和朝拜，都归于祢、圣父、圣子及圣神，从现在直到永远，永无穷尽。全体百姓应说：阿们。主教应说：\Quote{愿天主的平安与你们众人同在。}全体百姓应说：\Quote{也与你的心灵同在。}执事应再次宣告：——\par

% structure: p0033 source=h3 target=subsection reason=relative-heading-rank; base=h2

\subsection{神圣祭献后为信众的祈请祷词}

十三．让我们再进一步藉着祂的基督恳求天主，并因献给主天主的祭品而恳求祂，愿良善的天主藉着祂的基督的转求，在祂天上的祭台上接纳这祭品，作为馨香的祭品。让我们为这教会及民众祈祷。让我们为每一位主教、每一位司铎、所有在基督内的执事和服事者、为全体会众祈祷，愿主保守并保存他们众人。让我们为\BibleQuote{君王及一切有权位者}祈祷，使他们对我们和平相处，\BibleQuote{好使我们能敬虔端正、平安无事的度日}。\BibleRef{弟前 2:2} 让我们纪念圣殉道者，好使我们堪当分享他们的考验。让我们为那些在信德中安息的人祈祷。让我们为空气的调和及果实的完全成熟祈祷。让我们为新受光照的人祈祷，使他们能在信德中坚强，并使众人互相安慰。天主，求祢以祢的恩宠提拔我们。让我们站起来，并通过祂的基督将自己奉献于天主。然后主教说：天主，祢是伟大的，祢的名是伟大的，祢在计谋中是伟大的，在作为上是全能的，祢是祢的圣子耶稣、我们的救主的天主和父；求祢垂顾我们和祢的这群羊，就是祢藉着祂为光荣祢的名所拣选的；并圣化我们的身体和灵魂，恩赐我们得以\BibleQuote{洁净自己，除去身体和灵魂的一切污秽}，\BibleRef{格后 7:1} 并能获得为我们所存留的美物，不要认为我们中任何人是不配的；但求祢做我们的安慰者、帮助者和保护者，藉着祢的基督，愿光荣、尊威、赞美、颂谢和感谢归于祢和圣神，直到永远。阿们。此后，众人都说了阿们，执事说：请留意。然后主教向民众这样说：圣物是为圣洁的人。民众回答：只有一位是圣的；只有一位主，一位耶稣基督，永受赞美，归于天主圣父的光荣。阿们。\BibleQuote{荣归于至高天主，在地上平安，良善施于人。贺三纳于达味之子！奉主名而来的当受赞美，}祂就是向我们显现的主天主，\BibleQuote{贺三纳于至高天。}\BibleRef{路 2:14} 此后，主教领受，然后司铎、执事、副执事、读经员、咏唱者、苦修者；然后妇女中的女执事、贞女、寡妇；然后儿童；然后全体民众按次序，怀着敬畏和虔敬，不可喧哗。主教给予祭品，说：基督的身体；领受者说：阿们。执事拿起圣爵，递出时说：基督的血，生命的圣爵；饮用者说：阿们。在其余人领受时，咏唱第三十三篇圣咏。当所有男女都领受完毕，执事将剩余之物送入更衣所。咏唱者结束之后，执事说：——\par

% structure: p0035 source=h3 target=subsection reason=relative-heading-rank; base=h2

\subsection{领受后的求恩祈祷}

十四．现在我们已领受了基督宝贵的身体和宝贵的血，让我们感谢祂，因祂认为我们配得分享这些祂的神圣奥迹；并恳求祂，使我们领受这些不成为定罪的缘由，而是成为救恩，对灵魂和身体有益，保守虔敬，获得罪过的赦免，并得到来世的生命。让我们起来，并藉着基督的恩宠，将我们自己奉献于天主，于独一非受生的天主，及祂的基督。然后主教献上感恩：——\par

% structure: p0037 source=h3 target=subsection reason=relative-heading-rank; base=h2

\subsection{领受后的祈祷经文}

十五、全能的主天主，祢的基督、祢的受祝福的圣子的父，祢倾听那些以正直呼求祢的人，也知晓那些沉默者的恳求；我们感谢祢，因为祢认为我们配得分享祢的神圣奥迹，祢将这些奥迹赐予我们，为使我们完全确认我们所正确认识的一切，为保存虔诚，为赦免我们的过犯；因为祢的基督的名已呼求在我们身上，我们已与祢结合。祢将我们从恶人的共融中分别出来，求祢使我们与那些在圣洁中奉献给祢的人联合；以祢圣神的助佑，在真理中坚定我们；启示我们所不知道的，补足我们所欠缺的，坚定我们已经知道的；在祢的敬拜中保守司铎们无可指摘；使君王享有和平，使统治者秉行正义，使空气温煦适宜，使果实丰饶，以全能的天主眷顾治理世界；平息好战的国家，转化迷途的人，圣化祢的子民，保守守贞者，保存婚姻中的信德者，坚固纯洁者，使婴儿长大成人，坚定新领洗者；教导慕道者，并使他们堪当领受入门圣事；并藉我们的主耶稣基督，将我们全体聚集于祢的天国。愿荣耀、尊崇和敬拜，藉着祂，在圣神内，归于祢，至于无穷世。阿们。执事要说：藉着祂的基督向天主俯首，并领受祝福。主教要添加这篇祈祷，说：全能的天主，真天主，无与伦比者，祢无处不在，临于万物，却又不在任何事物中如同事物本身；祢不受地域限制，也不因时间而衰老；祢不被世代终结，也不被言语欺瞒；祢不受生育，无需护卫；祢超越一切朽坏，不受任何变化，本性不变；\BibleQuote{祢住在那不可接近的光中；}\BibleRef{弟前 6:16}祢本性不可见，却为所有以善意寻求祢的理性受造物所认识，并被那些以善意追求祢者所领悟；以色列的天主，祢的子民真正看见，并相信了基督：求祢恩待我，垂听我，为了祢圣名的缘故，祝福那些向祢俯首的人，赐予他们心中为他们的好处所求的恩惠，不要将其中任何人摒除于祢的王国之外；反要圣化、保护、遮蔽、助佑他们；拯救他们脱离那对头和一切仇敌；保守他们的家宅，护卫\BibleQuote{他们的出入。}因为荣耀、赞美、威严、敬拜和钦崇归于祢，并归于祢的圣子耶稣、祢的基督、我们的主、天主和君王，以及圣神，从今直到永远。阿们。执事要说：平安散去。这些关于这奥秘敬拜的宪章，我们宗徒们，为你们——主教、长老和执事——所制定。\par

% structure: p0039 source=h2 target=section reason=relative-heading-rank; base=h2

\section{第三节　圣职人员的祝圣与职责}

% structure: p0040 source=h3 target=subsection reason=relative-heading-rank; base=h2

\subsection{关于长老的祝圣——主所爱的\Person{若望}的宪章}

十六、关于长老的祝圣，我主所爱的\Person{若望}制定此宪章：主教啊，你应在长老团和执事面前，按手在他头上，祈祷说：全能的主、我们的天主，祢藉着基督创造了万物，同样藉着祂照料整个世界；因为祂有能力创造不同的受造物，也有能力按照它们不同的本性照料它们；因此，天主啊，祢藉着单纯的保存照料不朽的存在，藉着延续照料有死之物，藉着法律的提供照料灵魂，藉着需求的供给照料身体。求祢现在也垂顾祢的圣教会，使它增长，增多在其中的首领，赐予他们权能，使他们能以言语和行动为祢子民的建立而劳苦。现在也垂顾祢的这仆人，他由全体圣职人员的投票和决议被纳入长老团；求祢以恩宠和谋略之灵充满他，使他能以纯洁的心协助和治理祢的子民，如同祢曾垂顾祢的选民，命令\Person{梅瑟}拣选长老，并以祢的灵充满他们。上主，求祢现在也赐予这一切，并在我们内保存祢恩宠之神，使此人充满治病的恩赐和教导的言语，能以温和教导祢的子民，以纯洁的心和甘愿的灵魂忠诚地事奉祢，并为祢的子民完全履行神圣的职务。藉着祢的基督，愿荣耀、尊崇和敬拜归于祢和圣神，至于无穷世。阿们。\par

% structure: p0042 source=h3 target=subsection reason=relative-heading-rank; base=h2

\subsection{关于执事的祝圣——\Person{斐理伯}的宪章}

十七、关于执事的祝圣，我\Person{斐理伯}制定此宪章：主教啊，你应在全体长老团和执事面前，按手在他头上，并祈祷说：——\par

% structure: p0044 source=h3 target=subsection reason=relative-heading-rank; base=h2

\subsection{祝圣执事的祈祷格式。}

十八、全能的天主，真实而信实的天主，祢丰厚于一切以真理呼求祢的人，祢在谋划中可敬畏，在理解中有智慧，祢有大能且伟大。主啊，求祢垂听我们的祈祷，让祢的耳朵收纳我们的恳求，并\BibleQuote{使祢圣容的光辉照耀在祢的这仆人身上，}他是为祢被祝圣为执事职务的；求祢以祢的圣神和大能充满他，如同祢充满祢的殉道者、祢基督苦难的追随者\Person{斯德望}一样。使他堪当可接纳地履行执事的职务，从而藉此通过祢独生子的中介，得以达到更高的品位。愿荣耀、尊崇和敬拜归于祢和圣神，至于无穷世。阿们。\par

% structure: p0046 source=h3 target=subsection reason=relative-heading-rank; base=h2

\subsection{关于女执事——\Person{巴尔多禄茂}的宪章}

十九、关于女执事，我\Person{巴尔多禄茂}制定此宪章：主教啊，你应在长老团、执事和女执事面前，按手在她头上，并说：——\par

% structure: p0048 source=h3 target=subsection reason=relative-heading-rank; base=h2

\subsection{女执事祝圣祈祷文格式}

二十 永生的天主，我们主耶稣基督的圣父，造男造女的主，曾以圣神充满米黎盎、德波辣、亚纳和胡耳达；祢并未嫌弃祢的独生子由女人所生；祢也在会幕和圣殿中设立女子看守祢的圣门——求祢现在也垂顾这位将被祝圣为女执事的婢女，赐予她祢的圣神，并\BibleQuote{洁净她身灵的一切污秽}\BibleRef{格后 7:1}，使她相称地完成所托付的职事，以光荣祢，并赞美祢的基督，愿光荣与朝拜藉着祂归于祢和圣神，至于无穷。阿们。\par

% structure: p0050 source=h3 target=subsection reason=relative-heading-rank; base=h2

\subsection{论副执事——多默的约法}

二十一 关于副执事，我多默为你们主教制定此约法：主教啊，当你祝圣副执事时，你要按手在他身上，说：主天主，天地及其中万物的创造者，祢曾在会幕中指派监督和看守祢圣器的人；求祢现在也垂顾这位被立为副执事的仆人，赐予他圣神，使他能相称地处理祢职事的器皿，并永远奉行祢的旨意，藉着祢的基督，愿光荣、尊威和朝拜与祂一同归于祢和圣神，至于无穷。阿们。\par

% structure: p0052 source=h3 target=subsection reason=relative-heading-rank; base=h2

\subsection{论读经员——玛窦的约法}

二十二 关于读经员，我玛窦，又名肋未，曾为税吏，制定此约法：你要藉按手祝圣读经员，并向天主祈祷说：永生的天主，祢富于仁慈和怜悯，藉着祢的化工显明了世界的秩序，并保存祢选民数目；求祢现在也垂顾这位将被托付向祢百姓宣读祢圣经的仆人，赐予他祢的圣神，即先知之神。祢曾教导祢的仆人厄斯德拉向百姓宣读祢的法律，求祢现在也因我们的祈祷教导祢的仆人，并恩赐他无可指摘地完成所托付的职事，因而堪受更高的品级，藉着基督，愿光荣和朝拜与祂一同归于祢和圣神，至于无穷。阿们。\par

% structure: p0054 source=h3 target=subsection reason=relative-heading-rank; base=h2

\subsection{论宣信者——阿尔斐的儿子雅各伯的约法}

二十三 我阿尔斐的儿子雅各伯制定关于宣信者的约法：宣信者不需祝圣，因为他是由选择和忍耐而成的，并堪受极大的尊荣，因他曾在万民和君王面前宣认了天主和祂的基督的名。但若有需要，他可被祝圣为主教、司铎或执事。但若任何未受祝圣的宣信者因自己的宣认而擅自夺取这样的尊位，应剥夺并弃绝此人，因他并不在如此职位上，因为他否认了基督的约法，且\BibleQuote{比不信的人更坏}\BibleRef{弟前 5:8}。\par

% structure: p0056 source=h3 target=subsection reason=relative-heading-rank; base=h2

\subsection{同一宗徒论童贞女的约法}

二十四 我同一人制定关于童贞女的约法：童贞女不需祝圣，因为我们没有从主领受这样的命令\BibleRef{格前 7:25}；因为这完全是自愿试炼的状态，并非为贬损婚姻，而是为专务虔敬。\par

% structure: p0058 source=h3 target=subsection reason=relative-heading-rank; base=h2

\subsection{号称塔德乌的肋拜论寡妇的约法}

二十五 我号称塔德乌的肋拜制定关于寡妇的约法：寡妇不需祝圣；然而，若她丧夫已久，且生活庄重无可指摘，并特别照顾了家庭，如那些极负盛名的女子友弟德和亚纳，可被选入寡妇之列。但若她新近丧偶，则不可轻信，而应按时间判断她的青春；因为情感有时会随人年龄而衰老，若非有更好的约束加以抑制。\par

% structure: p0060 source=h3 target=subsection reason=relative-heading-rank; base=h2

\subsection{同一宗徒论驱魔者}

二十六 我同一人制定关于驱魔者的约法。驱魔者不需祝圣。因为这是自愿的行善尝试，并藉基督藉圣神的感动而得的恩宠。因为获得治愈之恩的人是由天主的启示而显明，他内在的恩宠对众人是明显的。但若有需要，他可被祝圣为主教、司铎或执事。\par

% structure: p0062 source=h3 target=subsection reason=relative-heading-rank; base=h2

\subsection{客纳罕人西满论祝圣主教所需的人数}

二十七 我客纳罕人西满制定约法，以确定主教应由多少人选举。应由三位或两位主教祝圣主教；但若仅由一位主教祝圣，则被祝圣者与祝圣者都应受剥夺。但若因迫害或类似原因，无法有更多主教聚集，而只得一位主教祝圣，则他必须取得较多主教的许可同意。\par

% structure: p0064 source=h3 target=subsection reason=relative-heading-rank; base=h2

\subsection{同一宗徒关于主教、司铎、执事及其他圣职人员的法规}

XXVIII. 有关宗徒法规，我同一位制定如下规定：主教祝福，但不接受祝福；他覆手、祝圣、奉献，接受主教们的祝福，但绝不受司铎的祝福。主教可剥夺任何应受剥夺的圣职人员，但主教除外，因为他自己没有权力这样做。司铎祝福，但不接受祝福；然而他从主教或同僚司铎接受祝福；同样地，他也向同僚司铎施予祝福。他覆手，但不祝圣；他不剥夺，但若属下应受此种惩罚，他可将他们隔离。执事不祝福，不给祝福，但从主教和司铎接受祝福；他不付洗，不奉献；但当主教或司铎奉献后，他分施给民众，不是作为司祭，而是作为服务司祭的人。其他圣职人员中，任何人不得擅自做执事的工作。女执事不祝福，也不做任何属于司铎或执事职务的事，只应看守门户，并在为妇女授洗时协助司铎，以保全端庄。执事可在司铎缺席时，隔离副执事、读经员、歌咏员或女执事（若有需要）。副执事、读经员、歌咏员或女执事不得隔离任何圣职人员或平信徒，因为他们是服事执事的人。\par

% structure: p0066 source=h2 target=section reason=relative-heading-rank; base=h2

\section{第四节　某些祈祷与法律}

% structure: p0067 source=h3 target=subsection reason=relative-heading-rank; base=h2

\subsection{关于祝福水与油——\Person{玛弟亚}的法规}

XXIX. 关于水与油，我\Person{玛弟亚}制定如下规定：主教应祝福水或油。若主教不在场，则由司铎祝福，执事侍立一旁。若主教在场，则司铎与执事侍立，主教如此祈祷：万军的上主，全能的天主，水的创造者，油的供给者，你是慈悲的、爱护人类的，你赐下水供人饮用和洁净，赐下油令人面容愉快喜乐；现今也请藉着你的基督，以奉献这些物品之人的名义，圣化这水与这油，并赐予它们恢复健康、驱除疾病、驱逐魔鬼、驱散一切陷阱的能力，藉着我们的希望基督，愿光荣、尊崇和敬拜归于你及圣神，直到永远。阿们。\par

% structure: p0069 source=h3 target=subsection reason=relative-heading-rank; base=h2

\subsection{同一位宗徒关于初果与什一之物的法规}

XXX. 我同一位制定关于初果与什一之物的规定：所有初果都应送至主教、司铎和执事，作为他们的生活所需；而所有什一之物则用于供养其余圣职人员、贞女、寡妇及处于贫困考验中的人。因为初果属于司铎及服事他们的执事。\par

% structure: p0071 source=h3 target=subsection reason=relative-heading-rank; base=h2

\subsection{同一位宗徒关于剩余祭品的法规}

XXXI. 我同一位制定关于剩余祭品的规定：在奥秘中剩余的祝福品，执事应按照主教或司铎的意向分给圣职人员：主教得四分，司铎得三分，执事得二分，其余副执事、读经员、歌咏员或女执事各得一分。因为这在主眼中是美好而可悦纳的，使每人按自己的位份受到尊敬；因为教会不是混乱的学校，而是良好秩序的学校。\par

% structure: p0073 source=h3 target=subsection reason=relative-heading-rank; base=h2

\subsection{\Person{保禄}宗徒关于自献领洗者的诸项法规——我们应接纳何人，应拒绝何人}

三十二. 我，\Person{保禄}，宗徒中最小的，也为你们——主教、司铎和执事——订立以下关于法典的宪典。那些首先来到敬虔奥秘的人，应由执事带到主教或司铎面前，并应查问他们为何来到主的圣言前；那些带领他们的人应仔细询问他们的品性，并为他们作证。应查问他们的品行和生活，以及他们是奴隶还是自由人。若有人是奴隶，应问其主人是谁。若他是一位信徒的奴隶，应问他的主人能否为他提供良好证明。若不能，应拒绝他，直到他向主人证明自己配得。但若主人给予良好证明，则应接纳他。但若他是一位外教人的家奴，应教导他讨主人喜欢，免得圣言被亵渎。那么，若他有妻子，或女子有丈夫，应教导他们彼此知足；但若他们未婚，应教导他们不要行淫，而要进入合法婚姻。但若他的主人是信徒，且知道他有淫行，却仍不给他娶妻，或给女子嫁夫，则应将他隔离；但若有人附魔，应教导他虔诚，但在被洁净之前不可领受共融；但若死亡临近，则可接纳。若有人是妓院主，他要么停止卖淫行径，要么被拒绝。若妓女前来，她要么停止卖淫，要么被拒绝。若雕刻偶像者前来，他要么停止其职业，要么被拒绝。若从戏院来的人，无论是男是女，或马车夫、或角斗士、或赛跑者、或奖品表演者、或奥林匹克竞技者、或在那些竞赛中吹笛、弹琴、弹竖琴者、或舞蹈教师、或小贩，要么停止其职业，要么被拒绝。若士兵前来，应教导他 \BibleQuote{不做不义之事，不诬告任何人，以所得的粮饷为满足} \BibleRef{路3:14}。若他服从这些规则，应接纳他；若拒绝，则被拒绝。犯有不可提及之罪的人——鸡奸者、娈童者、魔术师、巫术者、占星者、占卜者、使用咒语者、变戏法者、江湖郎中、制作护符者、施咒者、算命者、手相师、以及那些遇见你时观察鸟或猫的眼睛或脚、或噪音、或象征声音的人——这些人都应被考验一段时间，因为这类邪恶难以洗净；若他们停止那些行为，应接纳他们；若不答应，则被拒绝。一个不信主之人的妾，若只忠于其主人一人，应被接纳；但若她与他人不贞，则被拒绝。若信徒有妾，若她是婢女，他应停止那种方式，以合法方式结婚；若她是自由女子，他应以合法方式娶她；若不如此，则被拒绝。凡跟随外邦风俗或犹太传说者，要么改正，要么被拒绝。若有人跟随戏院的游戏、狩猎、赛马或格斗，要么停止，要么被拒绝。凡要成为慕道者的人，应作慕道者三年；但若有人勤奋，且对事业有善意，则应接纳他；因为所判断的不是时间长短，而是生活方式。教导者，即使是在平信徒中，若他熟悉圣言且品行庄重，就应教导；因为 \BibleQuote{他们都要受天主的教导} \BibleRef{若6:45}。所有信徒，无论男女，当他们早晨醒来，在工作之前，当洗净自己，祈祷；但若有任何要理讲授，信徒应优先考虑敬虔的圣言而非工作。信徒，无论男女，应仁慈对待仆人，正如我们在前面各卷中规定并在我们的书信中所教导的。\par

% structure: p0075 source=h3 target=subsection reason=relative-heading-rank; base=h2

\subsection{仆人不应在哪几日工作。}

三十三. 我\Person{伯多禄}和\Person{保禄}订立以下宪典。让奴隶工作五天；但在安息日与主的日子，让他们有空去教会接受敬虔的教导。我们说安息日是纪念创造，主的日子是纪念复活。让奴隶在圣周及其后一周休息——前一周纪念苦难，后一周纪念复活；他们需要受教谁受苦和复活，以及谁允许他受苦并使他复活。让他们在升天节休息，因为那是基督救世工程的完成。让他们在五旬节休息，因为圣神的降临，赐给那些信从基督的人。让他们在他的诞辰节休息，因为在那天，意外的恩宠赐给了人，即耶稣基督，天主的\OriginalTerm{圣言}{Logos}，应生于童贞玛利亚，为拯救世界。让他们在主显节休息，因为在那天，基督的天主性显现，因为圣父在洗礼中为他作证；并且\OriginalTerm{护慰者}{Paraclete}以鸽子的形状，向旁观者指出那被作证的那一位。让他们在宗徒的节日休息：因为他们被指派为你们的导师，带领你们归向基督，并使你们堪当领受圣神。让他们在首位殉道者\Person{斯德望}以及其他圣殉道者的节日休息，他们选择基督胜过自己的性命。\par

% structure: p0077 source=h3 target=subsection reason=relative-heading-rank; base=h2

\subsection{我们应在何时祈祷，以及为何祈祷。}

XXXIV. 你们要在清晨、第三时辰、第六时辰、第九时辰、傍晚和鸡鸣时奉献你们的祈祷：在清晨，感谢主赐给你们光明，带领你们度过黑夜，迎来白昼；在第三时辰，因为主在此时从比拉多处接受了定罪判决；在第六时辰，因为主在此时被钉十字架；在第九时辰，因为主被钉十字架时万物震动，仿佛因不虔敬的犹太人的狂妄企图而战栗，且不能忍受加于其主的伤害；在傍晚，感谢主赐给你们黑夜，以休息每日的劳作；在鸡鸣时，因为此时辰带来白昼来临的佳音，为适合光明的行动。但若因不信者而不能前往圣堂，你，主教，应在房屋里召集他们，使虔敬者不进入不虔敬者的集会。因为不是地方圣化人，而是人圣化地方。若不虔敬者占据了那地方，要避开，因为它被他们所亵渎。正如圣洁的司铎圣化一个地方，同样，亵渎者玷污它。若既不能在圣堂也不能在房屋里集会，则让每人独自歌唱、诵读和祈祷，或两三人一起。因为\BibleQuote{那里有两个或三个人因我的名聚在一起，我就在他们中间。}不要让一个信友与慕道者一起祈祷，即使在房屋里也不可；因为已获接纳者不应被未获接纳者玷污。不要让一个虔敬者与异端者一起祈祷，即使在房屋里也不可。因为\BibleQuote{光明与黑暗有什么相通呢？}凡与奴隶有关系的基督徒，无论男女，要么断绝关系，要么被开除。\par

% structure: p0079 source=h3 target=subsection reason=relative-heading-rank; base=h2

\subsection{\Person{雅各伯}，基督的弟兄，关于晚祷的宪章。}

XXXV. 我，\Person{雅各伯}，按肉身是基督的弟兄，却是\OriginalTerm{独生天主}{the only begotten God}的仆人，并由主自己和宗徒们被任命为\Place{耶路撒冷}的主教，兹规定如下：当傍晚时，你，主教，应召集教会；在点灯时重复咏唱圣咏后，执事应为慕道者、附魔者、已光照者和忏悔者祈祷，正如我们先前所说。但在这些人散去之后，执事应说：\Quote{所有信友，让我们向主祈祷。}在作了先前规定的祷词之后，他应说：——\par

% structure: p0081 source=h3 target=subsection reason=relative-heading-rank; base=h2

\subsection{晚祷的邀请祷词。}

XXXVI. 天主，求祢拯救我们，并藉祢的基督使我们兴起。让我们站立，恳求主的仁慈和祂的怜悯，为和平的天使，为美好和有益的事，为基督徒的离世，为平安且无罪的傍晚与黑夜；并让我们恳求我们一生的全部历程无可指摘。让我们将自己和彼此奉献给永生的天主，藉着祂的基督。然后，主教应加上这个祈祷，并说：——\par

% structure: p0083 source=h3 target=subsection reason=relative-heading-rank; base=h2

\subsection{晚祷的感恩。}

XXXVII. 天主，无始无终，藉基督创造整个世界，并予以照顾，但在万有之先，是祂的天主和父亲，圣神的主宰，理智与感官万物的君王；祢造了白日以行光明之事，造了黑夜以休息我们的软弱——因为\BibleQuote{白日属于祢，黑夜也属于祢；祢预备了光明和太阳。}——主，热爱人类者、诸善之源，求祢现在仁慈地接纳我们这晚间的感恩。祢带领我们度过了白日的长度，并将我们带到黑夜的开端，求祢藉着祢的基督保护我们，赐予我们平安的傍晚和无罪的黑夜，并藉着祢的基督恩赐我们永生，藉着祂，荣耀、尊威和朝拜，在圣神内归于祢，至于永远。阿们。执事应说：俯首接受覆手。主教应说：我们祖先的天主，仁慈的主，祢以祢的智慧造了人，为有理性的受造物，且比地上其他生物更受天主喜爱，并赐给他权柄管理地上的受造物，又按祢的旨意设立了统治者和司铎——前者为生活的安全保障，后者为有规律的敬拜——全能的主，求祢现在也垂顾，使祢的圣容光照祢的子民，他们低首俯心，求祢藉基督降福他们；藉着祂，祢以知识之光光照了我们，并向我们启示了祢自己；偕同祂，一切有理性和圣洁的本性都当向祢，并向那\OriginalTerm{护慰者}{the Comforter}圣神，献上合宜的朝拜，至于永远。阿们。执事应说：\Quote{平安离去。}同样，在早晨，在重复早晨的圣咏之后，在他遣散慕道者、附魔者、候洗者和忏悔者之后，在照常的邀请祈祷之后（以免我们重复相同的事），执事应在\Quote{天主，求祢拯救我们，并藉祢的恩宠使我们兴起}这句话之后加上：让我们恳求主的仁慈和祂的怜悯，使这个早晨和这一天平安无罪，也使我们整个旅程的时间如此；愿祂赐给我们祂的和平天使，基督徒的离世，并愿天主仁慈和宽厚。让我们将自己和彼此奉献给永生的天主，藉着祂的独生子。然后，主教应加上这个祈祷，并说：——\par

% structure: p0085 source=h3 target=subsection reason=relative-heading-rank; base=h2

\subsection{早晨的感恩。}

XXXVIII. 天主，神明和一切血肉的天主，无与伦比，一无所缺；祢赐太阳统治白日，月亮和星辰统治黑夜；求祢现在也以慈目垂顾我们，收纳我们早晨的感谢，并怜悯我们；因为我们没有\Quote{向陌生的天主伸开我们的双手；}我们中间没有新的天主，只有祢，永恒而无终的天主，祢藉基督赐给我们存在，又藉祂赐给我们福祉。也求祢藉祂赐给我们永生；愿光荣、尊威和敬拜归于祢和圣神，直到永远。阿们。执事说：低头接受按手。主教加上这祷文说：——\par

% structure: p0087 source=h3 target=subsection reason=relative-heading-rank; base=h2

\subsection{清晨的按手礼}

XXXIX. 天主，祢是信实真实的，祢向爱祢的人\Quote{广施仁慈于千万代，}喜爱谦卑的人，保护穷乏的人，万物都需要祢，因为万有都服从祢；求祢垂视这向祢低头的人民，以属灵的祝福祝福他们。\Quote{保护他们如同保护眼中的瞳仁，}保守他们虔敬公义，并藉着祢的爱子耶稣基督赐给他们永生。愿光荣、尊威和敬拜归于祢和圣神，从今直到永远。阿们。执事说：\Quote{平安离去。}当奉献初果时，主教如此感恩：——\par

% structure: p0089 source=h3 target=subsection reason=relative-heading-rank; base=h2

\subsection{初果的祈祷格式}

XL. 全能的主，全世界的创造者和保存者，我们藉着祢的独生子、我们的主耶稣基督，为奉献给祢的初果感谢祢，尽管不是我们应尽的方式，而是按我们的能力。因为谁能堪当地为祢赐予他们享用的一切而感谢祢呢？亚巴郎、依撒格、雅各伯和一切圣人的天主，祢藉祢的言语使万物结果实，命令大地生出各种果实，供我们喜悦和食物；祢赐给较迟钝和较温顺的生物汁液——给食草者青草，给一些动物肉，给另一些种子，但赐给我们谷物，作为有益且合适的食物，还有许多其他东西——有些为我们的需要，有些为我们的健康，有些为我们的愉悦。因此，为所有这些，祢藉基督堪受崇高的赞美诗歌，因为祢的恩惠。愿光荣、尊威和敬拜归于祢，在圣神内，直到永远。阿们。关于那些在基督内安息的人：在祈求祈祷之后，为避免重复，执事应补充如下：——\par

% structure: p0091 source=h3 target=subsection reason=relative-heading-rank; base=h2

\subsection{为亡者的祈求祷文}

XLI. 让我们为在基督内安息的弟兄祈祷，愿爱人的天主，收纳了他灵魂的，赦免他每一项罪，自愿的和非自愿的，并仁慈恩待他，赐他产业在虔诚者的地上，被送入亚巴郎、依撒格和雅各伯的怀中，与一切从创世以来中悦祂并遵行祂旨意的人同在，在那里一切忧愁、悲伤和哀哭都被驱逐。让我们起来，将我们自己和彼此奉献给永恒的天主，藉着那在起初就存在的圣言。主教说：祢本性不死，存在无穷尽，一切受造之物，无论不死或有死，都源于祢；祢造人成为理性受造物，这世界的公民，按其构成有死，又加给复活的应许；祢不许哈诺客和厄里亚尝到死亡的滋味：\BibleQuote{亚巴郎的天主，依撒格的天主，雅各伯的天主，祂不是死人的天主，而是活人的天主：因为所有人的灵魂都为祢而活，义人的灵魂在祢手中，没有痛苦能触及它们；}\BibleRef{玛 22:32}因为他们都在祢手中成圣。求祢现在也垂顾这位仆人，祢已拣选他并接纳他进入另一种状态，赦免他无论自愿或非自愿所犯的罪，赐给他仁慈的天使，安置他在圣祖、先知、宗徒和一切从创世以来中悦祢的人的怀中，那里没有忧伤、悲哀和哀哭，只有虔诚人的平安之地，义人不受搅扰的土地，以及在那里看见祢基督荣耀的人。愿光荣、尊威、敬拜、感谢和钦崇归于祢，在圣神内，直到永远。阿们。执事说：低头接受祝福。主教为他们感恩说：\Quote{主，求祢拯救祢的子民，祝福祢的产业，这是祢用祢基督的宝血所赎买的。在祢右手中牧养他们，在祢翅膀下荫庇他们，恩赐他们\InnerQuote{打那美好的仗，跑完赛程，保持信德}\BibleRef{弟后 4:7}坚定不移、无可指责、无可指摘，藉着我们的主耶稣基督，祢的爱子。愿光荣、尊威和敬拜归于祢和圣神，直到永远。阿们。}\par

% structure: p0093 source=h3 target=subsection reason=relative-heading-rank; base=h2

\subsection{我们应如何及何时庆祝已亡信友的纪念，并应当时常从他们的财产中分施给穷人。}

XLII. 亡者的第三日当以圣咏、读经和祈祷庆祝，因为祂在三日间复活了；第九日当为纪念活人和亡者而庆祝；第四十日依照古例：因为人民曾这样哀悼梅瑟，以及纪念他的周年日。应从他的财产中施舍给穷人，作为对他的纪念。\par

% structure: p0095 source=h3 target=subsection reason=relative-heading-rank; base=h2

\subsection{纪念或遗嘱对已亡的不敬神者毫无益处。}

XLIII. 我们这些话是对虔诚人说的；至于不敬神的人，即使你把全世界都给穷人，也丝毫不能帮助他。因为对谁而言神明在他活着时是仇敌，当他死后也必定仍是如此；因为天主没有不义。因为\Quote{上主是公义的，祂喜爱公义。}以及\BibleQuote{看哪，那人和他的工作。}\BibleRef{依 62:11}\par

% structure: p0097 source=h3 target=subsection reason=relative-heading-rank; base=h2

\subsection{论酗酒者。}

XLIV. 现在，当你们被邀请参加他们的纪念礼时，你们当以秩序和敬畏天主之情赴宴，准备为亡者代祷。因为你们是基督的司铎和执事，你们应当时常保持清醒，无论是在你们中间还是在外人面前，这样你们才能警告那些不守规矩的人。圣经说：\BibleQuote{掌权者有怒气。但他们不可饮酒，免得饮酒忘记了智慧，不能公正判断。}因此，司铎和执事在教会中仅次于全能天主和祂的爱子，拥有权威。我们说这些，并非要他们完全不饮酒，否则就是贬低天主为喜乐所造之物；而是不要因酒而迷乱。因为圣经并不说：不可饮酒；而是说什么？\BibleQuote{不可饮酒致醉；}又说：\BibleQuote{醉汉手中生出荆棘。}我们说这些不仅是对圣职人员，也是对每一位被称为我们主耶稣基督之名的平信徒基督徒。因为对他们也说过：\BibleQuote{谁有祸患？谁有忧愁？谁有争吵？谁有怨言？谁有红眼？谁无故受伤？这些不都是那些久坐饮酒、寻找酒席的人吗？} \BibleRef{箴 23:29}\par

% structure: p0099 source=h3 target=subsection reason=relative-heading-rank; base=h2

\subsection{论接待为基督受迫害者。}

XLV. 要接待那些因信仰受迫害、从这城逃到那城的人，\BibleRef{玛 10:23}要牢记主的话。因为知道虽然\BibleQuote{精神固然愿意，肉体却软弱，}\BibleRef{玛 26:41}他们就逃走，宁愿财物被夺，也要在自己身上保存基督的名而不否认。因此，供应他们所需，如此便成全了主的诫命。\par

% structure: p0101 source=h2 target=section reason=relative-heading-rank; base=h2

\section{第五节　所有宗徒敦促遵守教会的秩序}

% structure: p0102 source=h3 target=subsection reason=relative-heading-rank; base=h2

\subsection{每人当安于其位，不擅自攫取未交托之职务。}

第四十六章。现在，我们都共同命令你们，要每一个人都留在指派给他的品位中，不要超越自己的界限；因为这些界限不是我们的，而是天主的。因为主说：\BibleQuote{听你们的，就是听我；听我的，就是听派遣我来的那位。}并且说：\BibleQuote{轻视你们的，就是轻视我；轻视我的，就是轻视派遣我来的那位。}因为那些没有生命的事物也遵守良好的秩序，如同黑夜、白昼、太阳、月亮、星辰、元素、季节、月份、周、日、时辰，它们都服务于所指定的用途，正如所说的：\BibleQuote{你给他们设定了不可逾越的界限；}又关于海洋：\BibleQuote{我给它设定了界限，并用门闩和门将其围住；我对它说：你只可到这里，不可再往前；}\BibleRef{约 38:10}那么你们岂不更不应冒险改变我们按照天主旨意为你们所决定的事！但因为许多人认为这是小事，并冒险混淆品级，移动属于各自顺序的圣职授任，抢夺从未赐予他们的尊位，并以暴虐的方式擅自授予他们本身没有的权威，从而激起天主的愤怒（如同科辣黑的追随者和乌齐雅王所做的，他们没有权威，却未经天主的委任而篡夺了大司祭职；前者被火烧死，后者额头上长了癞病）；并且激起耶稣基督的愤怒，祂制定了这项法令；也令圣神忧伤，使祂的见证无效。因此，我们预知做这些事的人所面临的危险，以及由于不应献祭的人不敬地献祭而导致关于祭献和感恩祭的职务的疏忽，这些人认为大司祭职的尊荣（这是对我们君王耶稣基督这位大司祭的模仿）是儿戏；我们认为有必要就此事也提醒你们。因为有些人已经偏离到自己的虚妄中去了。我们说天主的仆人梅瑟（天主曾与他面对面说话，如同人同朋友说话；天主对他说：\BibleQuote{我在众人中认识你；}天主直接对他说话，而非藉着隐晦的方法、梦境、天使或谜语）——这个人当制定法令和神圣律法时，区分了哪些事应由大司祭执行，哪些应由司祭执行，哪些应由肋未人执行；将每个人在神圣礼仪中适当和合适的职务分配给他们。那些分配给大司祭做的事，司祭不得干涉；分配给司祭的事，肋未人不得干涉；每个人都遵守为他们所记载和指定的职分。若有任何人超越传统，其刑罚就是死亡。\Person{撒乌耳}的例子最清楚地表明了这一点，他以为可以在没有先知和大司祭\Person{撒慕尔}的情况下献祭，却为自己招来了无法补救的罪过和诅咒。连\Person{撒慕尔}曾傅油他为王这一事实也没有阻止先知（责备他）。但天主在\Person{乌齐雅}的事例中，\EditorialNote{编年纪下第二十六章}藉着更明显的效果显示了同样的事，祂毫不延迟地惩罚了这一越界行为，而那位疯狂贪求大司祭职的人也被从他的王国中剥夺了。至于我们中间发生的事，你们自己并非不知道。因为你们无疑知道，那些由我们称为主教、长老和执事的人，是藉着祈祷和覆手而立的；他们名称的差异显示了他们职务的差异。因为不是任何愿意的人都能被祝圣，正如\Person{雅洛贝罕}统治下那些虚假伪造的牛犊司祭的情况一样；\BibleRef{列上 13:33}只有被天主召唤的人才能。因为如果没有规则或品级的区分，用一个名称执行所有职务就足够了。但我们由主教导了事物的次序，将大司祭的职能分配给主教，司铎的职能分配给长老，以及在他们二者之下的服务分配给执事，好使神圣敬礼得以洁净地举行。因为执事不可献祭，不可付洗，也不可给予较大或较小的祝福。长老也不可施行圣职授任；因为扰乱这一秩序不符合神圣性。因为\BibleQuote{天主不是混乱的天主}，以免下级人员暴虐地僭取属于上级的职能，为自己制定新的法律体系，导致自己的祸害，不知道\BibleQuote{他们碰刺棒是困难的}；因为这样的人并非与我们或与主教们作战，而是与普世的主教和圣父的大司祭，我们的主耶稣基督作战。大司祭、司祭和肋未人是由天主最爱的\Person{梅瑟}祝圣的。我们十三位宗徒是由我们的救主祝圣的；而我\Person{雅各伯}和\Person{克肋孟}，以及同我们在一起的其他人，是由宗徒们祝圣的，以免我们再次列出所有主教的名单。并且，长老、执事、副执事和读经员，是由我们大家共同祝圣的。因此，那位本性是大司祭的，是独生子基督；祂没有自己夺取这尊荣，而是由圣父指定为这样的；祂为了我们而成为人，并向祂的天主和圣父献上了属神的祭献，在祂受难前独自将此任务交给我们执行，尽管当时还有其他人相信了祂。但相信的人不会立刻被任命为司铎，或获得大司祭的尊位。但在祂升天后，我们按照祂的制定，献上了纯洁无血的祭献；并祝圣了七位主教、长老和执事；其中一位是\Person{斯德望}，那位蒙福的殉道者，他在对天主的虔诚意向方面不逊于我们；他藉着对主耶稣基督的信德和爱德，显示了如此大的虔诚，以至于为他舍命，被杀害主的犹太人用石头砸死。然而，如此伟大而良善的人，充满圣神，看见基督坐在天主的右边，天国之门敞开，却从未见他行使不属于他执事职务的职能，也没有献祭，也没有为任何人覆手，而是将执事职分保持到底。因为这样才对得起这位为基督殉道的人，保持良好秩序。但如果有人责备我们的执事\Person{斐理伯}和我们的忠实兄弟\Person{阿纳尼雅}，因为前者为太监付洗，后者为我\Person{保禄}付洗，这些人并不明白我们所说的。因为我们只是断言没有人擅自夺取司铎的尊贵，而是要么从天主那里接受，如同\Person{默基瑟德}和\Person{约伯}，要么从大司祭那里接受，如同\Person{亚郎}从\Person{梅瑟}那里。因此，\Person{斐理伯}和\Person{阿纳尼雅}并没有自立，而是由基督——那位无可比拟的天主的大司祭——所指定的。\par

% structure: p0104 source=h3 target=subsection reason=relative-heading-rank; base=h2

\subsection{同一圣宗徒的教会法规。}

XLVII.1. 主教应由两三位主教祝圣。\par

2. 司铎由一位主教祝圣，执事及其余圣职人员亦然。\par

3. 若有主教或司铎，违背吾主关于祭献的规定，在祭台上奉献蜜、奶、烈啤代酒，或任何必需品，或鸟、兽、豆类等，超出规定的，应被剥夺圣职；但新谷粒、麦穗或当季的葡萄串除外。\par

4. 因为除了这些，以及在神圣奉献时的圣灯之油和乳香之外，不可在祭台上奉献任何东西。\par

5. 但所有其他果实应送到主教住所，作为给他和司铎们的初熟之果，而非送到祭台。显然，主教和司铎应将它们分给执事及其余圣职人员。\par

6. 主教、司铎或执事不得假借虔敬之名休弃自己的妻子；若休弃，应暂停神职。若继续如此，则剥夺圣职。\par

7. 主教、司铎或执事不得承担世俗事务；若承担，则剥夺圣职。\par

8. 若有主教、司铎或执事在春分前与犹太人一同庆祝逾越节，应剥夺圣职。\par

9. 若有主教、司铎、执事或司祭品册中的任何人，在奉献结束后不领受共融，应说明理由；若理由正当，可予宽恕；若不说明，应暂停神职，因他成为对人民的损害之由，并引起对奉献者的怀疑，仿佛奉献不当。\par

10. 凡信友进入圣教会，聆听圣经，但在祈祷和圣体共融时不留下者，应被暂停，因其在教会中造成混乱。\par

11. 若有人即使在自家房屋里与被绝罚者一同祈祷，他也应被暂停。\par

12. 若有圣职人员与被剥夺圣职者如同与圣职人员一同祈祷，他本人也应被剥夺圣职。\par

13. 若有圣职人员或平信徒，已被暂停或不应被接纳，却离开并在另一城市无推荐信而被接纳，则接纳者与被接纳者均应被暂停。若此人已被暂停，则延长其暂停期，因他向教会说谎欺骗。\par

14. 主教不应离开自己的教区跳往另一教区，即使群众强迫，除非有正当理由迫使他如此，例如他能以虔敬之言为新区人民带来更大益处；但这不应由他自己决定，而应由多位主教的裁决和非常恳切的请求。\par

15. 若有司铎、执事或圣职品册中的任何人，离开自己的教区前往另一教区，并完全迁居，在未经其主教同意的情况下继续留在该教区，我们命令他不得继续其职务，尤其当其主教召他返回而他却不服从，继续处于混乱状态。但他在那里可以平信徒身份领受共融。\par

16. 但若接纳他们的主教轻视对他们所定的剥夺令，接纳他们为圣职人员，则应暂停该主教，作为混乱的导师。\par

17. 受洗后结过两次婚，或有姘妇者，不能被立为主教、司铎、执事或司祭品册中的任何人。\par

18. 娶了寡妇、被休的妇人、娼妓、奴婢或戏子者，不能被立为主教、司铎、执事或司祭品册中的任何人。\par

19. 娶了两姐妹，或娶了侄女、甥女者，不能成为圣职人员。\par

20. 圣职人员若成为担保人，应被剥夺圣职。\par

21. 阉人，若因遭受人的伤害或在迫害中失去\Emph{生殖器}，或天生如此，且堪当主教职者，可被立为主教。\par

22. 自残者不得被立为圣职人员，因他是自杀者，是反对天主创造的敌人。\par

23. 圣职人员若自残，应被剥夺圣职，因他是自杀者。\par

24. 平信徒若自残，应被隔绝三年，因他自设陷阱。\par

25. 主教、司铎或执事若犯奸淫、伪证或偷窃，应被剥夺圣职，但不可暂停神职；因为圣经说：\BibleQuote{你不可因同一罪行两次报复而加苦害。}\par

26. 对余下的圣职人员，同样处理。\par

27. 对于以单身身份加入圣职的人，我们只准许读经员和歌咏员，若他们愿意，以后可以结婚。\par

28. 我们命令：主教、司铎或执事若击打犯罪的信友或作恶的不信者，以为可以借此恐吓他们，应被剥夺圣职，因为吾主从未教导我们这样的事。反之，\Quote{当祂自己被击打时，祂没有还手；祂被辱骂时，没有还口；祂受苦时，没有恐吓。}\par

29. 若有主教、司铎或执事，因明显的罪行被公正地剥夺圣职，竟敢干预曾托付给他的职务，应将该人完全逐出教会。\par

30. 若有主教用金钱获得该职位，或司铎、执事亦然，则该人与祝圣他者均应被剥夺圣职，并完全断绝共融，如同西满术士被我伯多禄所行。\par

31. 若有主教利用世俗统治者，并藉其手段获得某教会的主教职位，应剥夺其圣职并暂停神职，凡与他共融者亦然。\par

32. 若有司铎轻视自己的主教，并另外聚集，另立祭台，而他在主教身上无可指摘其在虔敬或正义方面的缺失，则此人应作为野心家被剥夺圣职，因他是暴君，凡依附他的其余圣职人员亦然。平信徒应被暂停。但这些事须在受到主教一次、两次甚至三次劝诫之后方可执行。\par

33. 若有司铎或执事被其主教停职，其他任何人不得接纳他，唯有将他停职的那位主教可以，除非将他停职的主教去世。\par

34. 不可接纳任何陌生的主教、司铎或执事，除非持有推荐信；若有推荐信，应予以审核。若他们是虔诚的宣讲者，应被接纳；若非如此，则供应其所需，但不可接纳他们共融：因为许多事是出于疏忽而做的。\par

35. 各地的主教应当知道他们中间谁是首位，并尊他为他们的首领，未经他同意，不可做任何大事；但每人只管理属于自己教区及其下属地方的事务。然而，未得全体同意，他不可做任何事；因为只有这样，才能有同心合意，天主也因基督在圣神内得光荣。\par

36. 主教不可冒险在自己辖区之外，为不属于自己的城市或地区行圣秩授予。若他被证实未经那些城市或地区管理者同意而如此行，则他本人及其所祝圣者都应被剥夺圣职。\par

37. 若有已祝圣的主教不承担其职务，也不照顾交托给他的百姓，应被停职，直至他承担为止；司铎和执事亦然。但若他前往就职，却因民众的顽固态度而非因他本人不愿而未被接纳，则应继续担任主教；但该城的圣职人员应被停职，因为他们没有更好地教导那不顺服的百姓。\par

38. 主教会议应每年举行两次，彼此询问虔诚的教义，并解决发生的教会争议——一次在五旬节期的第四周，另一次在\OriginalTerm{希珀伯雷泰奥斯月}{Hyperberetæus}十二日。\par

39. 主教应负责管理教会财物，如同在天主面前一样管理。但他不可将任何部分占为己有，也不可将属于天主之物给予自己的亲属。若亲属贫穷，他应像对待穷人一样供养他们；但不可假借这些名义挪用教会收入。\par

40. 司铎和执事不可未经主教同意做任何事，因为主教受托管理主的百姓，并须为他们的灵魂交账。主教自己的财产（若有）与属于主的财产应公开区分，使他去世时能按照自己的意愿将财产留给他想给的人；以免假借教会收入之名，使主教自己的财产受损——他有时有妻子儿女、亲属或仆人。因为这在天主和人面前是公正的：既不让教会因不知哪些是主教自己的收入而受损失，也不让他的亲属假借教会之名而陷入困境，或使他的亲人陷入诉讼，以致他的去世受人指责。\par

41. 我们命令主教对教会财物拥有管理权；因为如果他受托管理人的珍贵灵魂，他更应当对财物发出指示，以便根据他的权威，由司铎和执事将所有财物分配给有需要的人，并怀着对天主的敬畏和所有敬意用于他们的供养。他也可以取用他所需的那些东西，如果他确实需要，以供他自己和与他同住的弟兄们的必要之用，使他们绝不陷入窘迫：因为天主的法律规定了在祭台侍奉的人应当由祭台供给；因为连士兵也从不自备军饷与敌人作战。\par

42. 凡沉溺于掷骰子或酗酒的主教、司铎或执事，要么停止这些行为，要么被剥夺圣职。\par

43. 若副执事、读经员或歌咏员也有类似行为，要么停止，要么被停职；平信徒亦同。\par

44. 凡向借贷者索取利息的主教、司铎或执事，要么停止这样做，要么被剥夺圣职。\par

45. 只与异端者同祷的主教、司铎或执事，应被停职；但若他还容许异端者履行圣职的任何部分，应被剥夺圣职。\par

46. 我们命令，凡接受异端者的洗礼或祭献的主教、司铎或执事，应被剥夺圣职：\BibleQuote{基督与贝利亚尔有什么相合？或者信主的人与不信主的人有什么相干？}\par

47. 若主教或司铎重新给已领受真洗礼者施洗，或不给被不虔敬者玷污的人施洗，应被剥夺圣职，因为他嘲笑主的十字架和死亡，不能分辨真正的司祭和假冒者。\par

48. 若平信徒休弃自己的妻子而另娶，或娶被他人休弃的妻子，应被停职。\par

49. 若任何主教或司铎不按主的定规施洗，即因圣父、圣子、圣神之名，而施洗入三位无始者、或三子、或三位护慰者，应被剥夺圣职。\par

50. 若任何主教或司铎不执行一次入门的三次浸水，而只执行一次浸水（即归于基督之死），应被剥夺圣职；因为主并没有说：\BibleQuote{给我施洗归于我的死}，而是说：\BibleQuote{你们要去使万民成为门徒，因圣父、圣子、圣神之名给他们施洗。}所以，诸位主教，要按照基督的旨意和我们的圣神所定的规条，三次施洗入一位圣父、圣子、圣神。\par

51. 若任何主教、司铎、执事，或任何在圣职册上的人，出于憎恶而禁绝婚姻、肉和酒，并非为了自身操练，且忘记\BibleQuote{天主所造的一切都非常好}\BibleRef{创 1:31}，以及\BibleQuote{天主照自己的肖像造了男人和女人}\BibleRef{创 1:26}，并亵渎地辱骂受造物，要么他改正，要么被剥夺圣职并被逐出教会；平信徒亦同。\par

52. 若有主教或司铎不接受那从罪恶中悔改归来的人，反而拒绝他，就应停职；因为他使基督忧伤，基督说：\BibleQuote{对于一个悔改的罪人，在天上也有喜乐。} \BibleRef{路 15:7}\par

53. 若有主教、司铎或执事在庆节不食肉或不饮酒，就应停职，因为他有如\BibleQuote{良心已经烙上印记}，且成为许多人的绊脚石。 \BibleRef{弟前 4:2}\par

54. 若有圣职人员被发现在酒馆中进食，就应停职，除非他不得已而在旅途客栈中歇脚。\par

55. 若有圣职人员无端辱骂自己的主教，就应停职；因为圣经说：\BibleQuote{不可诽谤你人民的官长。} \BibleRef{出 22:28}\par

56. 若有圣职人员辱骂司铎或执事，就应隔离。\par

57. 若有圣职人员嘲笑瘸子、聋子、瞎子或跛足的人，就应停职；平信徒亦同。\par

58. 若有主教或司铎不照顾圣职人员和民众，也不以虔诚教导他们，就应隔离；若他继续疏忽，就应褫夺圣职。\par

59. 若有主教或司铎，当某位圣职人员匮乏时，不供给其急需，就应停职；若他继续如此，就应褫夺圣职，如同杀害了自己的弟兄。\par

60. 若有人在教会公开诵读不敬虔之人的伪书，仿佛它们是圣书，致使民众和圣职人员败坏，就应褫夺圣职。\par

61. 若有基督徒被控犯奸淫、通奸或任何其他禁行，且被定罪，则不得晋升至圣职。\par

62. 若有圣职人员因畏惧人，如畏惧犹太人或外邦人或异端者，而否认基督之名，就应停职；若他否认圣职人员的名分，就应褫夺圣职；但当他悔改时，应视作平信徒接纳。\par

63. 若有主教、司铎、执事，或任何列于司祭名册者，吃带有生命之血的肉，或被野兽撕裂的肉，或自死的肉，就应褫夺圣职；因为法律本身禁止了此事。若他是平信徒，就应停职。\par

64. 若有圣职人员被发现于主日或安息日（除仅有一次外）守斋，就应褫夺圣职；若他是平信徒，就应停职。\par

65. 若任何圣职人员或平信徒进入犹太人或异端者的会堂祈祷，就应褫夺圣职并停职。\par

66. 若有圣职人员因争吵击打某人，且一击致死，就应因其鲁莽而褫夺圣职；若他是平信徒，就应停职。\par

67. 若有人对未许配的童女施暴并留住她，就应停职。但不可另娶他人；必须保留他所选定的女子，即使她贫穷。\par

68. 若有主教、司铎或执事从任何人接受再次祝圣，就应褫夺其圣职，祝圣他的人亦然，除非他能证明先前的祝圣来自异端者；因为由这类人所受洗或祝圣者，既不能成为基督徒，也不能成为圣职人员。\par

69. 若有主教、司铎、执事、读经员或唱经员，不守四十日大斋，或不守每周第四日（周三）和预备日（周五）的斋戒，就应褫夺圣职，除非因身体软弱受阻。若他是平信徒，就应停职。\par

70. 若有主教或其他圣职人员与犹太人一同守斋，或与他们一同过庆节，或接受他们庆节的礼物，如无酵饼或类似物品，就应褫夺圣职；若他是平信徒，就应停职。\par

71. 若有基督徒带油进入异教神庙，或进入犹太人的会堂，或在他们的庆节点灯，就应停职。\par

72. 若任何人，无论是圣职人员或平信徒，从圣教会取走蜂蜜或油，就应停职，并应偿还所取之物外加五分之一。\par

73. 凡祝圣过的银器、金器或麻布器皿，任何人不得据为己用，因这是不义的；若有人被捉住，应受停职处罚。\par

74. 若有主教被可信赖的忠实人士控告犯任何罪行，其他主教必须传唤他；若他前来并为自己辩护，但仍被定罪，则规定其处罚。但若传唤时他不服从，应第二次派遣两位主教去传唤他。若他仍藐视他们，拒绝前来，则主教会议可对他作出任何其认为适当的判决，以免他因逃避审判而显得得利。\par

75. 不接受异端者作证指控主教；也不接受独身的基督徒。因为法律说：\BibleQuote{凭两三个证人，一切事可成立。}\par

76. 主教不得以主教尊位满足其兄弟、儿子或任何其他亲属，或随意祝圣任何人；因为将主教职作为继承，并在神圣之事上顺遂人情是不义的。我们不可将天主的教会置于继承法之下；若有人如此行，其祝圣无效，并应受停职处罚。\par

77. 若有人一只眼伤残，或腿瘸，但配得主教尊位，应立为主教；因为玷污他的不是身体的缺陷，而是灵魂的污染。\par

78. 但若他既聋又瞎，则不可立为主教；并非因他是受玷污之人，而是为避免妨碍教会事务。\par

79. 若有人附有魔鬼，不可立为圣职人员。甚至不得与信众一同祈祷；但当他被洁净后，应接纳他；若他配得，则可受祝圣。\par

80. 不可立即祝圣刚由外邦人归化并受洗的人为主教，或来自邪恶生活方式的人；因为尚未经受任何考验的人就成为他人的导师，这是不义的，除非在某些情况下出于天主的恩宠。\par

81. 我们说，主教不应让自己卷入公共行政事务，而应随时专心处理教会的必要事务。因此，要么他同意不这样做，要么被剥夺职务。因为，按照主的训诫，\Quote{没有人能事奉两个主人，} \BibleRef{玛 6:24}。\par

82. 我们不允许在没有主人同意的情况下将仆人授予神职；因为这会伤害他们的主人。因为这种做法会导致家庭的颠覆。但如果有时一个仆人看起来配得被授予高级神职，就像我们\OriginalTerm{敖乃息摩}{Onēsimus}那样，并且他的主人允许，给予他自由，并让他离开自己的家，那么让他被授予神职。\par

83. 凡加入军队，并希望同时保留罗马政府和司祭职务的主教、司铎或执事，应被剥夺职务。因为\Quote{凯撒的归凯撒，天主的归天主。}\par

84. 凡无理辱骂君王或总督者，应受惩罚；若是神职人员，应被剥夺职务；若是平信徒，则应被停职。\par

85. 以下书籍，你们无论是神职人员还是平信徒，都应当视为可敬和神圣的。旧约：梅瑟五书——\BibleBook{}{创世纪}、\BibleBook{}{出谷纪}、\BibleBook{}{肋未纪}、\BibleBook{}{户籍纪}、\BibleBook{}{申命纪}；\BibleBook{若苏厄}{书}一卷，\BibleBook{}{民长纪}一卷，\BibleBook{}{卢德传}一卷，\BibleBook{}{列王纪}四卷，\BibleBook{}{编年纪}两卷，\BibleBook{}{厄斯德拉}两卷，\BibleBook{}{艾斯德尔传}一卷，\BibleBook{}{友弟德传}一卷，\BibleBook{}{玛加伯}三卷，\BibleBook{}{约伯传}一卷，\BibleBook{}{圣咏}一百五十篇；撒罗满三书——\BibleBook{}{箴言}、\BibleBook{}{训道篇}、\BibleBook{}{雅歌}；先知书十六卷。此外，务必让你们的年轻人学习博学的\OriginalTerm{息辣}{Sirach}的智慧。至于我们的圣书，即新约，则是：\BibleBook{玛窦}{福音}、\BibleBook{马尔谷}{福音}、\BibleBook{路加}{福音}和\BibleBook{若望}{福音}四福音；保禄十四封书信；伯多禄两封书信，若望三封，雅各伯一封，犹达一封；\Person{克肋孟}两封书信；以及奉献给你们主教，由我\Person{克肋孟}所著的\OriginalTerm{宗徒宪章}{Constitutions}，共八卷；因其包含奥秘，不宜在所有人面前公开；还有我们宗徒的\BibleBook{宗徒}{大事录}。\par

让这些教会法规由我们为你们，诸位主教，确立下来；如果你们持续遵守，你们将得救，并享有和平；但如果你们不顺服，你们将受到惩罚，彼此之间永无休止的战争，并承受与你们不顺服相称的惩罚。\par

现在，唯独非受生的天主，整个世界的创造者，借着祂的平安，在圣神内将你们众人联合；使你们在一切善工上得以完全，坚定不移，无可指责，无可挑剔；并经由祂的爱子、我们的天主和救主耶稣基督的调解，赐予你们与我们同在的永生；偕同祂，愿光荣归于你，万有的天主和圣父，在护慰者圣神内，从今时直到永远。阿们。\par

由\Person{克肋孟}所著的\WorkTitle{圣宗徒宪章}（即公教教义）终。\par

\SourceRecord{Translated by James Donaldson.；Ante-Nicene Fathers；Vol. 7.；Edited by Alexander Roberts, James Donaldson, and A. Cleveland Coxe.；Buffalo, NY: Christian Literature Publishing Co.,；1886.；<http://www.newadvent.org/fathers/07158.htm>.}
